% !TeX program = xelatex
% ================================================================
%  末尾论文《非线性波动方程的几何与傅里叶方法》独立重排
%  —— 独立文档，复用 mathtranslation.cls v3.1 的字体、定理环境与中文排版
%
%  层级对应: 每“讲”为一个 \chapter，讲内小节为 \section*。
%  (mathtranslation.cls v3.1 基于 ctexbook，顶层为 \chapter；
%   原文每“讲”即本文的最高层级，故由旧 ctexart 版的 \section 上移为 \chapter。)
% ================================================================
\documentclass[Biblatex=false]{mathtranslation}
\usepackage{enumitem}

% ---------- 定理类环境：按“讲”（章）共享 X.N 计数器 ----------
% mathtranslation.cls 默认各环境独立计数；本论文原书将命题/定理/引理/定义/注
% 统一编号为“讲号.序号”，故在此解除原定义并重建。
% 注意：只 \let 环境命令，不释放 \c@theorem 等计数器本身。
\makeatletter
\let\theorem\relax\let\endtheorem\relax
\let\definition\relax\let\enddefinition\relax
\let\lemma\relax\let\endlemma\relax
\let\corollary\relax\let\endcorollary\relax
\let\proposition\relax\let\endproposition\relax
\let\remark\relax\let\endremark\relax
\let\example\relax\let\endexample\relax
\let\problem\relax\let\endproblem\relax
\makeatother

% 论文将命题/定理/引理/定义/注/问题/猜想统一按“讲号.序号”编号，
% 故新建共享计数器 mtthm（按“讲”即章重置）；“例”用独立计数器 mtx。
\newcounter{mtthm}[chapter]
\renewcommand{\themtthm}{\arabic{chapter}.\arabic{mtthm}}
\newcounter{mtx}[chapter]
\renewcommand{\themtx}{\arabic{chapter}.\arabic{mtx}}

\theoremstyle{mtfang}
\newtheorem{theorem}[mtthm]{定理}
\newtheorem{definition}[mtthm]{定义}
\newtheorem{lemma}[mtthm]{引理}
\newtheorem{corollary}[mtthm]{推论}
\newtheorem{proposition}[mtthm]{命题}
\newtheorem{remark}[mtthm]{注}
\theoremstyle{mtsong}
\newtheorem{example}[mtx]{例}
\newtheorem{problem}[mtthm]{问题}
\newtheorem{conjecture}[mtthm]{猜想}

% 无编号环境（用于不编号的“定理”“问题”“猜想”等块）。
% “例”“注”的菱形结束标记由类文件统一提供，此处不再重挂（重挂会出现双菱形）；
% 仅为新定义的带 * 未编号环境补挂，保持与编号版本外观一致。
\theoremstyle{mtfang}
\newtheorem*{theorem*}{定理}
\newtheorem*{definition*}{定义}
\newtheorem*{lemma*}{引理}
\newtheorem*{corollary*}{推论}
\newtheorem*{proposition*}{命题}
\newtheorem*{remark*}{注}
\newtheorem*{problem*}{问题}
\theoremstyle{mtsong}
\newtheorem*{example*}{例}
\newtheorem*{conjecture*}{猜想}
\AtEndEnvironment{example*}{\envenddiamond}
\AtEndEnvironment{remark*}{\envenddiamond}

% ---------- 公式编号：全书连续，与原书一致 ----------
% ctexbook（book 基类）默认把 equation 按章编号（\theequation = 章号.序号）。
% 本论文原书公式连续编号 (1)--(170)，且正文中大量引用（如"方程 (16)""(67)"）
% 为硬编码文本，故解除按章重置、恢复连续编号。
\makeatletter
\@removefromreset{equation}{chapter}
\renewcommand{\theequation}{\arabic{equation}}
\@removefromreset{figure}{chapter}
\renewcommand{\thefigure}{\arabic{figure}}
\makeatother

% ---------- 按类型引用的命令 ----------
\newcommand{\thmref}[1]{定理~\ref{#1}}
\newcommand{\lemref}[1]{引理~\ref{#1}}
\newcommand{\defref}[1]{定义~\ref{#1}}
\newcommand{\propref}[1]{命题~\ref{#1}}
\newcommand{\corref}[1]{推论~\ref{#1}}
\newcommand{\remref}[1]{注~\ref{#1}}
\newcommand{\exref}[1]{例~\ref{#1}}

% ---------- 标题块（极简，不套用书籍封面）----------
\title{非线性波动方程的几何与傅里叶方法}
\author{Sergiu Klainerman\thanks{普林斯顿大学}\\[4pt] \small 译}
\date{}
\makeatletter
\renewcommand{\maketitle}{%
  \begin{center}
    {\LARGE\headfont \@title}\\[1.6em]
    {\large \@author}%
  \end{center}%
  \vspace{1.2em}%
  \@thanks
}
\makeatother

% book 类无 abstract 环境，补一个极简版（居中“摘要”+ 缩进小字）。
\newenvironment{abstract}{%
  \par\medskip
  \begin{center}{\headfont 摘\quad 要}\end{center}%
  \begin{quote}\small\noindent
}{%
  \end{quote}\medskip
}

\begin{document}

% ================================================================
%                       前置：标题、摘要、目录
% ================================================================
\frontmatter
\maketitle

\begin{abstract}
非线性波动方程这一课题，在过去的二十年间因三项根本性进展而发生了深刻变化：
一、精密几何技巧的出现，即只涉及物理空间中几何构造的方法；
二、Littlewood--Paley 分解的系统运用以及仿微分法的引入；
三、傅里叶时空技巧的发展，如 Strichartz 型不等式与双线性估计。
本讲稿旨在讨论这些进展中若干技术性问题，并通过它们所导出的重要新结果加以说明。
我们也计划讨论若干最新进展，它们要求几何与傅里叶方法之间高度非平凡的合作，并预示着这两种看似不可调和的观点未来强有力的融合。
\end{abstract}

\newpage
\begingroup
\pagestyle{empty}
\tableofcontents
\clearpage
\endgroup

% ================================================================
%                         正文
% ================================================================
\mainmatter

\chapter{第一讲\quad 标准波动方程的估计}

考虑 Minkowski 时空 $\mathbb{R}^{n+1}$ 中的标准波动方程
\begin{equation}
\Box\phi=0.\end{equation}

Minkowski 空间 $\mathbb{R}^{n+1}$ 中的标准惯性坐标记为 $x^\mu$, $\mu = 0,1,\ldots,n$, 相对这组坐标 Minkowski 度量取对角形式 $m_{\mu\nu}= \mathrm{diag}(-1,1,\ldots,1)$. 我们有 $x^0 = t$, $x = (x^1,\ldots,x^n)$ 表示空间坐标. 下文中我们始终沿用关于重复指标求和的约定, 以及关于向量与张量指标升降的标准约定. 特别地, 若 $x_\mu = m_{\mu\nu}x^\nu$, 则 $x_0 = -t$, $x_i = x^i$, $i = 1,\ldots,n$. 我们采用记号 $\partial_\mu = \frac{\partial}{\partial x^\mu}$ 与 $\partial_t = \frac{\partial}{\partial x^\mu}$(原文如此). 记 $\Sigma_{t_0}$ 为类空超平面 $t = t_0$. 波动算子定义为 $\Box = m^{\alpha\beta}\partial_{\alpha\beta} = -\partial_t^2 + \sum_i\partial_i^2$.
下面我们回顾方程 (1) 在 $t=0$ 处满足初值问题的解的基本已知估计:
\begin{equation}
\phi(0,x)=f(x),\qquad \partial_{t}\phi(0,x)=g(x).\end{equation}

为方便起见, 记 $\phi[0] = (f, D^{-1}g)$, 其中 $D^{-1}$ 为象征 $|\xi|^{-1}$ 的拟微分算子.
\begin{proposition}[能量恒等式、不等式]\label{prop:1.1}
方程 (1)、(2) 的解满足
\begin{equation}
E[\partial\phi](t)=E[\partial\phi](0),\end{equation}
其中
\begin{equation}
E[\partial\phi](t)=\int_{\Sigma_{t}}\bigl(|\partial_{t}\phi|^{2}+\sum_{i}|\partial_{i}\phi|^{2}\bigr)\,d x.\end{equation}
由此可得能量不等式: 对任意 $s \geq 0$,
\begin{equation}
\|\partial\phi(t)\|_{H^{s}(\mathbb{R}^{n})}\leq\|\partial\phi(0)\|_{H^{s}(\mathbb{R}^{n})}.\end{equation}
\end{proposition}

\begin{proof}
能量恒等式既可由只涉及分部积分的几何方法证明, 也可利用 Plancherel 公式与傅里叶表示公式, 通过傅里叶方法证明:
\begin{equation}
\phi(t,x)=(2\pi)^{-n}\int e^{i x\cdot\xi}\Bigl(\cos t|\xi|\,\hat f(\xi)+\frac{\sin t|\xi|}{|\xi|}\,\hat g(\xi)\Bigr)\,d\xi.\end{equation}
更高阶的能量不等式基于 $\Box$ 与 $\partial_{\mu}$ 的可交换性.\end{proof}

\begin{remark}\label{rem:1.2}
标准的 Sobolev 嵌入 $H^s(\mathbb{R}^n) \subset L^\infty(\mathbb{R}^n)$(当 $s > \frac{n}{2}$)使我们无需利用显式表示即可得到方程 (1) 解的 $L^\infty$ 估计. 这一手续在与局部存在定理的联系中可以推广到非线性方程.\end{remark}

\begin{proposition}[色散不等式]\label{prop:1.3}
方程 (1)、(2) 的解满足
\begin{equation}
\|\phi(t)\|_{L^{\infty}}\leq c\,t^{-\frac{n-1}{2}}\|D^{\frac{n+1}{2}}\phi[0]\|_{L^{1}}.\end{equation}
\end{proposition}

\begin{remark}\label{rem:1.4}
事实上 (6) 并不十分准确; 如果把左边的 $L^\infty$ 范数换成 $\mathrm{BMO}$ 范数, 或者把右边的 $L^1$ 范数换成 Hardy 范数 $\mathcal{H}^1$, 则正确的估计成立. 不过, 若数据 $\phi(0) = f$、$\partial_t \phi(0) = g$ 的傅里叶变换支在某个二进球壳 $\frac{\lambda}{2} \leq |\xi| \leq 2\lambda$(其中某个固定的 $\lambda \in 2^{\mathbb{N}}$)中, 则 (6) 如所书写是成立的.\end{remark}

\begin{proof}
直到最近, (6) 唯一稳健的证明都基于对表示式 (5) 的驻相法. 在奇数维, 利用解的角均值表示也可以证明色散估计的一种相关形式. 我们后面将讨论一种避免任何表示公式的 (6) 的推导.\end{proof}

\begin{remark}\label{rem:1.5}
色散不等式提供了两类信息. 第一类涉及当 $t \to \infty$ 时 $\|\phi(t)\|_{L^\infty}$ 的精确衰减率; 第二类涉及对 $t > 0$ 时 $\|\phi(t)\|_{L^\infty}$ 的正则性性质. 就改进的正则性而言, 与 Sobolev 嵌入 $L^\infty(\mathbb{R}^n) \subset W^{1,n}(\mathbb{R}^n)$ 相比, (6) 在 $t > 0$ 时多获得了 $\frac{n-1}{2}$ 阶导数.\end{remark}

众所周知, 就渐近行为而言, (6) 在应用于非线性波动方程时并不十分有用. 推导波动方程解渐近性质的一个更有效的方法, 基于由交换向量场法得到的广义能量估计, 以及全局 Sobolev 不等式. 下面我们极快地介绍这一观点, 并在第一讲的补充中更系统地展开. 我们回顾用交换向量场法推导上述衰减率估计. 其思想是将能量恒等式 (3) 与交换向量场以及经典 Sobolev 不等式的全局形式结合使用.
Minkowski 时空 $\mathbb{R}^{n+1}$ 装备有一族 Killing 与保形 Killing 向量场: 平移 $T_{\mu} = \partial_{\mu}$, Lorentz 旋转 $O_{\mu\nu} = x_{\mu} \partial_{\nu} - x_{\nu} \partial_{\mu}$, 伸缩 $S = t \partial_{t} + x^{i} \partial_{i}$, 以及反平移 $K_{\mu} = -2x_{\mu}S + \langle x, x\rangle \partial_{\mu}$. 这里 $x^{\mu}$ 表示标准变量 $x^{0} = t, x^{1}, \ldots, x^{n}$, 而 $x_{\mu} = m_{\mu\nu} x^{\mu}$. Killing 向量场 $T_{\mu}$ 与 $O_{\mu\nu}$ 与 $\Box$ 交换, 而 $S$ 以保持解空间的意义与 $\Box$ 交换: 由 $[\Box, S] = 2\Box$ 知 $\Box \phi = 0$ 蕴含 $\Box S \phi = 0$. 我们把算子 $O_{\mu\nu}$ 分解为角旋转算子 $(ij)$: $O = x_{i} \partial_{j} - x_{j} \partial_{i}$ 与推进 $(i)$: $L = x_{i} \partial_{t} + t \partial_{i}$, 其中 $i, j, k = 1, \ldots, n$. 回顾 (4) 中的能量表达式; 基于上述交换性质我们定义如下"广义能量"
\begin{equation}
E_{k}[\partial\phi]=\sum_{X_{i_{1}},\ldots,X_{i_{j}}}E^{2}\bigl[\partial X_{i_{1}}X_{i_{2}}\ldots X_{i_{j}}\phi\bigr],\end{equation}
其中求和取遍 $0 \leq j \leq k$ 以及所有 Killing 向量场 $T, L_{\mu\nu}$ 与伸缩向量场 $S$. 交换向量场法的要点在于, 量 $E_k\;(k \geq 1)$ 被方程 (1) 的解守恒. 因此, 若对所有 $0 \leq k \leq s$ 数据 $f, g$ 满足
\begin{equation}
\int(1+|x|)^{2k}\Bigl(|\nabla^{k+1}f(x)|^{2}+|\nabla^{k} g(x)|^{2}\Bigr)\,d x\leq C_{s}<\infty,\end{equation}
则对所有 $t$ 有 $E_s[\partial\phi](t) \leq C_s$. 方程 (1) 的所需衰减估计现在可由下面的 Sobolev 不等式的全局形式(见 \cite{klai27})推出:
\begin{proposition}[衰减估计, 见 {\cite{klai27}}]\label{prop:1.6}
设 $\phi$ 为 $\mathbb{R}^{n+1}$ 中任一函数, 使得对某整数 $s > \frac{n}{2} + 1$ 有 $E_s[\phi]$ 有限. 则对所有 $t > 0$,
\begin{equation}
|\partial\phi(t,x)|\lesssim(1+t+|x|)^{-\frac{n-1}{2}}(1+|t-|x||)^{-\frac{1}{2}}E_{s}[u].\end{equation}
因此, 若 (1) 中的数据 $f, g$ 对 $0 \leq k \leq s$ 满足 (8) 且某 $s > \frac{n}{2}$, 则对所有 $t \geq 0$,
\begin{equation}
\|\partial\phi(t,\cdot)\|_{L^{\infty}}\lesssim\frac{1}{(1+t+|x|)^{\frac{n-1}{2}}(1+|t-|x||)^{\frac{1}{2}}}.\end{equation}
\end{proposition}

\begin{remark}\label{rem:1.7}
显然, 这个证明纯为几何$^{1}$ 的估计蕴含了色散不等式 (6) 给出的衰减性质. 事实上, 它在波区 $|x| \sim t$ 之外提供了更多信息, 与线性方程 $\Box \phi = 0$ 预期的传播性质十分吻合. 另一方面, 由于 (9) 实质上是 Sobolev 不等式的全局形式, 似乎\propref{prop:1.6} 的估计与 (6) 改进的正则性特征无关. 然而这并不正确, 我们将在下一个命题中看到.\end{remark}

\begin{remark}\label{rem:1.8}
交换向量场法(\cite{klai27}, 亦见 \cite{klai46})最初是作为 \cite{klai17} 几乎整体存在结果的另一种大为简化的证明而发展的. 该方法非常灵活, 适用于证明维数 $n>3$ 的整体存在性以及维数 $n=3$ 对一大类非线性波动方程(包括完全非线性的)在小初值下的几乎整体存在性. 它还与零条件(见\propref{prop:1.22} 后的讨论)结合, 证明低维 $n=3,2$ 的整体存在性. 向量场法还被许多其他作者发展, 特别是 H. Lindblad \cite{klai48} 与 T. Sideris \cite{klai63,klai64,klai65}, 亦见 Hörmander \cite{klai16} 及其中的参考文献. 向量场法的一个重要修正出现在 \cite{klai12} 中, 其中向量场必须被修改. 这将在第三讲中一定程度上讨论.\end{remark}

\begin{proposition}[见 {\cite{klai30}}]\label{prop:1.9}
交换向量场法蕴含色散不等式 (6).\end{proposition}

\begin{proof}
不失一般性可设 $\partial_t \phi = g = 0$, 且 $f = \phi(0)$ 的傅里叶变换支在某个壳 $\frac{\lambda}{2} \leq |\xi| \leq 2\lambda$($\lambda \in 2^{\mathbb{N}}$)中. 通过简单的伸缩论证, 实际上可设 $\lambda = 1$. 由于 $\phi$ 关于空间变量 $x$ 的傅里叶变换 $\hat{\phi}$ 也支在同一壳中, 只需对 $\nabla \phi$ 或 $\nabla^k \phi$ 证明估计即可. 接下来用并排半径为有限、互不重叠的盘 $D_I$ 覆盖 $\mathbb{R}^n$, 这些盘以整数坐标点 $I \in \mathbf{Z}^n$ 为中心, 使得每个 $D_I$ 至多与有限个 $c_n$ 个盘 $D_J$ 相交, 其中 $c_n$ 只依赖于维数 $n$. 考虑光滑单位分解 $(\chi_I)_{I\in\mathbf{Z}^n}$, 满足 $\operatorname{supp}\chi_I\subset D_I$ 且各 $\chi_I$ 为正. 显然可以安排使得对所有 $k$,
\[
\sum_{I\in\mathbf{Z}^{n}}|\nabla^{k}\chi_{I}(x)|\leq C_{k,n}
\]
对 $x \in \mathbb{R}^n$ 一致成立. 现在令 $f_I = \chi_I \cdot f$, $\phi_I$ 为对应的方程 (1) 的解, 数据 $\phi_I(0) = f_I, \partial_t \phi_I(0) = 0$. 显然 $f = \sum_I f_I$, $\phi = \sum_I \phi_I$. 只需证明对所有 $I$,
\[
\|\nabla^{k}\phi_{I}(t)\|_{L^{\infty}}\leq C_{n,k}(1+t)^{-\frac{n-1}{2}}\sum_{j=0}^{n+k+1}\|D^{j}f_{I}\|_{L^{1}},
\]
其中常数 $C_{n,k}$ 只依赖于 $n$ 与 $k$. 确实, 若 (12) 成立, 则易得 $\|\nabla^{k}\phi(t)\|_{L^\infty} \leq C_{n,k}(1+t)^{-\frac{n-1}{2}} \sum_{j=0}^{n+k+1} \|\sum_{I} \nabla^{j}\chi_{I}\|_{L^\infty} \|f\|_{L^{1}}$, 从而由 (11) 得 $\|\nabla^{k}\phi(t)\|_{L^\infty} \leq C_{n,k}(1+t)^{-\frac{n-1}{2}} \|f\|_{L^{1}}$.
于是只需验证 (12). 经空间平移不妨设 $I = 0$. 把\propref{prop:1.6} 用于 $\psi = \nabla \phi_0$, 对 $s_*$ 为严格大于 $\frac{n}{2} + 1$ 的首个整数, 得
\[
\begin{aligned}
\|\psi(t)\|_{L^{\infty}}&\leq c(1+t)^{-\frac{n-1}{2}}E_{s_{*}}[\phi_{0}](t)\\
&\leq c(1+t)^{-\frac{n-1}{2}}E_{s_{*}}[\phi_{0}](0).\end{aligned}
\]
由于 $\phi_{0}$ 的支集包含于以原点为中心、半径为 1 的球内, 有
\[
E_{s_{*}}[\phi_{0}](0)\leq C_{n}\sum_{j=0}^{s_{*}+1}\|D^{j}f_{0}\|_{L^{2}}.
\]
最后, 根据 $\mathbb{R}^n$ 中的标准 Sobolev 不等式 $\|f\|_{L^2} \leq c\|\nabla^{\frac{n}{2}}f\|_{L^1}$, 得到
\[
\|\psi(t)\|_{L^{\infty}}\leq c(1+t)^{-\frac{n-1}{2}}\sum_{j=0}^{n+2+1}\|D^{j}f_{0}\|_{L^{1}},
\]
如所求.\end{proof}

在能量不等式与色散不等式之间作插值, 得到所谓的 Strichartz--Brenner 结果:
\begin{equation}
\|\phi(t)\|_{L^{r}}\leq c\,|t|^{-\gamma(r)}\|\nabla^{\sigma}\partial\phi(0)\|_{L^{r^{\prime}}},\end{equation}
其中 $\gamma(r) = (n-1)\bigl(\frac{1}{2}-\frac{2}{r}\bigr)$, $\frac{1}{r}+\frac{1}{r'}=1$, $r\geq2$, 且伸缩条件为 $\frac{n}{r}=-\gamma(r)-\sigma-1+\frac{n}{r'}$. 通过标准的 $TT^{*}$ 论证、Hardy--Littlewood--Sobolev 不等式以及 Littlewood--Paley 理论的应用, 得到广义 Strichartz 型不等式.
\begin{definition}\label{def:1.10}
一对正实数 $(q, r)$ 称为波容许(wave admissible), 如果
\[
\frac{2}{q}\leq\gamma(r)=(n-1)\Bigl(\frac{1}{2}-\frac{2}{r}\Bigr),\qquad q\geq2,\qquad(q,r,n)\neq(\infty,1,3).
\]
\end{definition}

\begin{proposition}[Strichartz 型不等式]\label{prop:1.11}
初值问题 (1)--(2) 的解满足
\[
\|\phi\|_{L_{t}^{q}L_{x}^{r}([0,T]\times\mathbb{R}^{n})}
\leq c\,\|\phi[0]\|_{H^{\sigma}},\qquad \sigma=n\Bigl(\frac{1}{2}-\frac{2}{r}\Bigr)-\frac{1}{q}.
\]
\end{proposition}

\begin{remark}\label{rem:1.12}
Strichartz 本人只证明了各向同性情形 $q = r$ 的估计. 其余估计归于 Pecher、Ginibre--Velo、Lindblad--Sogge 与 Keel--Tao. 关于这些估计的最终版本与精确参考文献, 我们参阅 Keel--Tao \cite{klai20}.\end{remark}

\begin{remark}\label{rem:1.13}
Strichartz 型不等式在非线性波动方程理论的许多最新进展中起着关键作用. 注意, 在固定频率下由能量恒等式与色散不等式推出 (13) 的步骤相当柔和, 它们可追溯到 Duhamel 原理与初值问题的唯一性$^{2}$. 两者都适用于一般变系数线性波动方程, 且只要求系数有很少的正则性. 因此 Strichartz 型不等式的主要构件是 (3) 与 (6).\end{remark}

迄今为止讨论的所有估计都只涉及波动方程 (1) 的单个解. 它们在应用于非线性项含导数的波动方程时价值有限. 过去十年中发展出一种新观点: 直接估计涉及 (1) 的多个解的二次甚至多次线性表达式, 可获得远为灵活的估计.
\begin{proposition}[双线性 Strichartz, 见 {\cite{klai47}}]\label{prop:1.14}
设 $\phi, \psi$ 为 (1) 的两个解. 设 $n \geq 2$ 且 $(q, r)$ 为波容许. 则
\[
\|D^{\sigma}(\phi\psi)\|_{L_{t}^{q/2}L_{x}^{r/2}}\lesssim\|\phi[0]\|_{\dot{H}^{s_{1}}}\|\psi[0]\|_{\dot{H}^{s_{2}}},
\]
其中 $\sigma, s_{1}, s_{2}$ 满足
\[
0<\sigma<n-\frac{2n}{r}-\frac{4}{q},\qquad
0<s_{1},s_{2}<\frac{n}{2}-\frac{n}{r}-\frac{1}{q},\qquad
s_{1}+s_{2}+\sigma=n-\frac{2n}{r}-\frac{2}{q}.
\]
\end{proposition}

\begin{proof}
证明基于如下色散不等式的微局部形式:\end{proof}

\begin{proposition}[改进的色散不等式, 见 {\cite{klai47}}]\label{prop:1.15}
考虑 (1)(2) 的一个解, 其初数据的傅里叶变换支在半径为小 $\mu$ 的球内, 该球包含于单位二进环 $\frac{1}{2} \leq |\xi| \leq 2$ 中. 则
\[
|\phi(t)|_{L^{\infty}}\leq c\,\mu\, t^{-\frac{n-1}{2}}\|\phi[0]\|_{L^{1}}.
\]
\end{proposition}

改进的色散不等式易由标准驻相法得到. 我们至今还没有满意的几何推导.
现在考虑更一般的双线性估计, 形如
\[
\left\|D^{\gamma}D_{+}^{\gamma+}D_{-}^{\gamma-}\left(\phi\psi\right)\right\|_{L_{t}^{q/2}L_{x}^{r/2}}\lesssim\|\phi[0]\|_{\dot{H}^{s_{1}}}\|\psi[0]\|_{\dot{H}^{s_{2}}},
\]
其中 $\phi$ 与 $\psi$ 满足 (1). 这里 $D^\gamma$、$D_{+}^{\gamma+}$、$D_{-}^{\gamma-}$ 为象征 $|\xi|^\gamma$、$|\tau| + |\xi||^{\gamma+}$ 与 $|\tau| - |\xi||^{\gamma-}$ 的乘子算子. 在 $q, r = 4$ 的情形, 所有此类估计都已知道. 下述定理的特例先后出现在 \cite{klai31}、以及 \cite{klai37,klai44,klai39} 中. 完整解最近由 Foschi--Klainerman \cite{klai14}(亦见 \cite{klai69})给出.
\begin{proposition}[双线性估计, 见 {\cite{klai14}}]\label{prop:1.16}
设 $n \geq 2$ 且 $\gamma, \gamma_{-}, \gamma_{+}, s_1, s_2 \in \mathbb{R}$. 估计
\[
\left\|D^{\gamma}D_{+}^{\gamma+}D_{-}^{\gamma-}\left(\phi\psi\right)\right\|_{L_{t}^{2}L_{x}^{2}}\lesssim\left\|\phi[0]\right\|_{\dot{H}^{s_{1}}}\left\|\psi[0]\right\|_{\dot{H}^{s_{2}}}
\]
对 (1) 的所有初数据成立, 当且仅当下列条件满足:
\[
\gamma+\gamma_{+}+\gamma_{-}=s_{1}+s_{2}-\frac{n-1}{2},\qquad
\gamma_{-}\geq-\frac{n-3}{4},\qquad
\gamma>-\frac{n-1}{2},
\]
\[
s_{i}\leq\gamma_{-}+\frac{n-1}{2}\quad (i=1,2),\qquad
s_{1}+s_{2}\geq\frac{1}{2},
\]
\[
\left(s_{i},\gamma_{-}\right)\neq\left(\frac{n+1}{4},-\frac{n-3}{4}\right)\quad (i=1,2),\qquad
\left(s_{1}+s_{2},\gamma_{-}\right)\neq\left(\frac{1}{2},-\frac{n-3}{4}\right).
\]
\end{proposition}

这些估计与零二次型密切相关.
\begin{definition}\label{def:1.17}
设 $\phi, \psi$ 为 $\mathbb{R}^{n+1}$ 上的两个光滑函数. 定义
\[
Q_{0}(\phi,\psi)=m^{\alpha\beta}\partial_{\alpha}\phi\partial_{\beta}\psi,\qquad
Q_{i j}(\phi,\psi)=\partial_{i}\phi\partial_{j}\psi-\partial_{i}\psi\partial_{j}\phi,\qquad
Q_{0i}(\phi,\psi)=\partial_{i}\phi\partial_{t}\psi-\partial_{i}\psi\partial_{t}\phi.
\]
\end{definition}

\begin{definition}\label{def:1.18}
若 $\phi, \psi$ 为 $\mathcal{S}'(\mathbb{R}^{n+1})$ 中的分布, 其时空傅里叶变换 $\widehat{\phi}, \widehat{\psi}$ 为函数. 记 $\phi \preceq \psi$ 当且仅当 $|\widehat{\phi}| \leq C\widehat{\psi}$.\end{definition}

\begin{lemma}[见 {\cite{klai44}}]\label{lem:1.19}
下列成立:
\[
Q_{0}(\phi,\psi)\precsim D_{+}D_{-}(\phi\psi),\qquad
Q_{i j}(\phi,\psi)\prec{\cal D}^{\frac{1}{2}}{\cal D}_{-}^{\frac{1}{2}}({\cal D}^{\frac{1}{2}}\phi\,{\cal D}^{\frac{1}{2}}\psi).
\]
\end{lemma}

\begin{remark}\label{rem:1.20}
对 $Q_{0i}$ 也可推出类似公式. 可结合\propref{prop:1.16} 的双线性估计与上面的引理, 导出零二次型的非平凡 $L^{2}$ 时空估计.\end{remark}

\begin{remark}\label{rem:1.21}
尽管\propref{prop:1.16} 的双线性估计中正则性的获得是近十年的进展, 但零二次型对非线性方程应用的重要性此前已在衰减获得(gain of decay)的联系中为人所知. 下面是一个典型事实.\end{remark}

\begin{proposition}[零形式的衰减, 见 {\cite{klai28}}]\label{prop:1.22}
设 $\phi,\psi$ 为 (1) 的两个解, 并设其初数据对充分大的 $s$ 满足加权 Sobolev 范数估计 (8). 则对任意零形式 $Q_{0},Q_{ij},Q_{0i}$, 有衰减估计
\[
\|Q(\phi,\psi)(t)\|_{L^{2}(\mathbb{R}^{n})}\leq C\, t^{-n}
\]
(当 $t \to \infty$).\end{proposition}

证明是\propref{prop:1.6} 衰减估计以下改进的直接推论.
\begin{proposition}[剥离衰减性质]\label{prop:1.23}
设 $\phi$ 为 $R^{n+1}$ 中使得对某整数 $s > \frac{n}{2} + 2$ 有 $E_s[\phi]$ 有限的函数. 令 $E_\pm = \partial_r \pm \partial_r$(其中 $\partial_r = \sum_i \frac{x^i}{|x|}\partial_i$), 且 $E_i = 0$, $E_A, E_B \geq \delta_{AB}$. 记 $|\nabla \phi|^2 = \sum_{i=1}^{n-1} E_A(\phi)^2$. 我们有
\[
\|E_{+}\phi(t)\|_{L^{\infty}}\lesssim\left(\frac{1}{t}\right)^{\frac{n+1}{2}},
\]
\[
\|\bar{\nabla}\phi(t)\|_{L^{\infty}}\lesssim\left(\frac{1}{t}\right)^{\frac{n+1}{2}},
\]
\[
\|\left(1+|u|\right)E_{-}\phi(t)\|_{L^{\infty}}\lesssim\left(\frac{1}{t}\right)^{\frac{n-1}{2}}.
\]
\end{proposition}

\begin{remark}\label{rem:1.24}
\propref{prop:1.22} 中零二次型改进衰减的应用, 见 \cite{klai28}, 用于证明 $\mathbb{R}^{3+1}$ 中一大类满足零条件的拟线性波动方程的小初值整体存在性. 证明 \cite{klai28} 中结果所用的向量场技巧被 Sideris 在 \cite{klai63}、\cite{klai64}、\cite{klai65} 中推广到 Wave Maps 与非线性弹性力学.\end{remark}

至此我们只考虑了 $q = r = 4$ 型 (16) 的一般双线性估计. 对于"波容许"指数 $q, r$ 与 $\gamma_{-}$, 当 $s_{1} = s_{2}$ 时\propref{prop:1.14} 给出的估计是最优的. 存在一个一般双线性猜想, 断言在指数 $\gamma, \gamma_{+}, \gamma_{-}, s_1, s_2$ 的适当限制下估计 (16) 成立. T. Tao \cite{klai68} 利用 T. Wolff \cite{klai76} 在其锥限制猜想近期重要工作中发展的思想, 证明了该猜想的一些非常重要情形. 不过, 我想以这一近期进展中一个重要而非常简单的要素来结束第一讲. 这就是一般波动可分解为管(tubes)的思想.
\begin{definition}\label{def:1.25}
称波 $\phi$ 为 $\pm$ 波, 如果其初数据满足
\[
\partial_{t}\phi(0)=\pm i\sqrt{-\Delta}\,\phi(0),
\]
或在傅里叶表示中 $\partial_t\hat\phi(\theta,\xi) = \pm i|\xi|\hat\phi(\theta,\xi)$. 注意任意波可按公式
\[
\phi=\phi^{+}+\phi^{-}
\]
分解为 $\pm$ 波, 其中 $\phi^{\pm}$ 的初数据为 $\phi^{\pm}[0] = P^{\pm}\phi[0]$,
\[
P_{\pm}\phi[0]:=\Bigl(\frac{\phi(0)\mp i(-\Delta)^{-1/2}\phi_{t}(0)}{2},\frac{\phi_{t}(0)\pm i(-\Delta)^{1/2}\phi(0)}{2}\Bigr).
\]
\end{definition}

\begin{remark}\label{rem:1.26}
注意解 $\phi^+, \phi^-$ 由其数据唯一确定. 换言之, 为定义它们只需分解 $\phi$ 的数据. 不过它们可用傅里叶变换表示为
\[
\phi^{\pm}(t,x)=\int e^{\pm i t|\xi|}e^{i x\cdot\xi}\widehat{P^{\pm}\phi[0]}(\xi)\,d\xi.
\]
\end{remark}

\begin{proposition}\label{prop:1.27}
考虑 $\phi = \phi^{-}$ 为一支减波, 其 $\phi^{-}[0]$(本质上)支在球 $B = B(0, \Lambda_0)$ 中, $\Lambda_0 > 1$, 且其傅里叶变换支在区域 $\{\xi: \frac{1}{2} \leq \xi_1 \leq 2,\ |\xi'| \leq \frac{1}{\Lambda_0}\}$ 中, 其中 $|\xi'| = \sqrt{\xi_2^2 + \cdots + \xi_n^2}$. 对任意 $t \geq 0$ 考虑区域 $T^+(t, \Lambda) = \{x: |x| \leq t + \Lambda_0;\ |x_1 - t|^2 + |x'|^2 \leq \Lambda^2\}$. 则对所有满足 $\Lambda_0 \leq \Lambda \leq t \leq \Lambda_0^2$ 的 $t$:
\[
\|\partial\phi^{-}(t)\|_{L^{2}(\Sigma_{t}\backslash T^+(t,\Lambda))}\leq C_{N}\Bigl(\frac{\Lambda_{0}}{\Lambda}\Bigr)^{N}E(\phi)^{\frac{1}{2}}.
\]
\end{proposition}

\begin{proof}
该命题的标准证明基于积分表示公式 (20) 中的简单分部积分. 也可给出基于交换向量场 \cite{klai43} 的纯几何证明.\end{proof}

上述命题的向量场法证明, 结合\propref{prop:1.22} 的零估计以及由 T. Wolff \cite{klai76} 与 T. Tao \cite{klai68} 首创的"尺度归纳"(induction on scales)重要思想, 使我们$^{4}$ 能给出如下双线性估计的纯几何证明:
\begin{proposition}\label{prop:1.28}
设 $\phi, \psi$ 为 (1) 的两个解, $Q$ 为零形式 $Q_{0}, Q_{ij}$ 或 $Q_{0i}$ 之一. 则
\[
\|Q(\phi,\psi)\|_{L^{2}(\mathbb{R}^{n+1})}\lesssim\|D\phi[0]\|_{L^{2}(\mathbb{R}^{n})}\|D^{\frac{n+1}{2}}\phi[0]\|_{L^{2}(\mathbb{R}^{n})}.
\]
\end{proposition}

这一估计是\propref{prop:1.16} 的特例, 最早由 Machedon 与我和 \cite{klai31} 借助 $Q(\phi, \psi)$ 的时空傅里叶变换大量依赖地证明. 它已被用于证明 Yang--Mills \cite{klai34}、Maxwell--Klein--Gordon \cite{klai33} 与临界幂 Yang--Mills--Higgs \cite{klai19} 的有限能量整体存在性.

\chapter{第二讲\quad 半线性波动方程的应用}

我们关注重要几何系统(如波动映射、Yang--Mills 与 Maxwell--Klein--Gordon 方程)的柯西问题. 作为第一近似, 所有这些方程都可写成形式
\begin{equation}
\Box\phi=\mathcal{N}(\phi),\end{equation}
其中 $\Box = -\partial_t^2 + \Delta$ 是 $\mathbb{R} \times \mathbb{R}^n$ 上的标准波动算子, $\Delta = \sum_1^n \partial_j^2$ 是 $\mathbb{R}^n$ 上的 Laplacian, $\phi = \phi(t, x)$ 取值于某 $N \geq 1$ 的 $\mathbb{R}^N$, 而 $\mathcal{N}$ 是关于时间变量局部的非线性算子(即 $\mathcal{N}(\phi)(t, \cdot)$ 只依赖于 $\phi(t', \cdot)$, 其中 $t'$ 在 $t$ 的任意邻域中).
我们在初值超平面 $t = 0$ 上给定柯西数据:
\begin{equation}
(\phi,\partial_{t}\phi)_{t=0}=(f,g)\in H^{s}\times H^{s-1},\end{equation}
其中 $H^s = \{f : (I - \Delta)^{s/2} f \in L^2\}$. 为简便记 $\phi[0] \in H^s$.
我们事实上集中关注以下类型的系统:
\section*{1. 波动映射(Wave Maps)}

\begin{equation}
\Box\phi^{I}+\sum_{J,K}\Gamma_{J K}^{I}(\phi)Q_{0}(\phi^{J},\phi^{K})=0.\end{equation}

这里 $\phi^{I}$ 表示 $\phi$ 的第 $I$ 个分量函数, $\Gamma_{JK}^{I}$ 是从 $\mathbb{R}^{N}$ 到 $\mathbb{R}$ 的光滑函数, $Q_{0}$ 为零形式
\[
Q_{0}(\phi,\psi)=\sum_{\mu=0}^{n}\partial_{\mu}\phi\partial^{\mu}\psi=-\partial_{t}\phi\partial_{t}\psi+\sum_{j=1}^{n}\partial_{j}\phi\partial_{j}\psi.
\]

\section*{2. Yang--Mills 型}

\begin{equation}
\Box\phi=D^{-1}Q(\phi,\phi)+Q(D^{-1}\phi,\phi),\end{equation}
其中 $D^\alpha = (-\Delta)^{\alpha/2}$, $Q$ 表示如下类型的任一双线性算子: 给定向量值函数 $\phi$ 与 $\psi$, $Q(\phi, \psi)$ 的第 $I$ 个分量函数是对所有 $1 \leq i < j \leq n$ 与所有 $J, K$ 的 $Q_{ij}(\phi^J, \psi^K)$ 的常实系数线性组合, 其中 $Q_{ij}$ 为零形式
\[
Q_{i j}(\phi,\psi)=\partial_{i}\phi\partial_{j}\psi-\partial_{i}\psi\partial_{j}\phi.
\]
((YM) 右端的两个 $Q$ 可能代表两个不同的此类算子. )

\section*{3. Maxwell--Klein--Gordon 型}

\begin{equation}
\left\{\begin{aligned}\Box\phi&=D^{-1}Q(\psi,\psi),\\ \Box\psi&=Q(D^{-1}\phi,\psi),\end{aligned}\right.\end{equation}
其中 $\phi = (\phi^1, \ldots, \phi^{N_1})$, $\psi = (\psi^1, \ldots, \psi^{N_2})$, $N = N_1 + N_2$, 而 $Q$ 意义同前. 因此 (MKG) 是 (YM) 的特例.
下面的定理概括了当今已知的主要适定性结果.
\begin{theorem*}[主定理]
\begin{enumerate}[label=\roman*.]
\item (\cite{klai35,klai44}. )若 $n \geq 2$ 且 $s > s_c = \frac{n}{2}$, 则 (WM) 对 $H^s(\mathbb{R}^n)$ 中的初值局部适定.
\item (\cite{klai72}、\cite{klai73}. )对 $n \geq 2$, 系统 (WM) 在 Besov 空间 $B_{2}^{\frac{n}{2},1}(\mathbb{R}^{n})$ 中对小数据局部适定.
\item (\cite{klai70}、\cite{klai41}、\cite{klai71}. )若 $n \geq 5$ 且目标流形承认有界可平行化结构, 则波动映射系统在 $H^{\frac{n}{2}}(\mathbb{R}^{n})$ 中对小数据"弱适定". 若目标流形为标准球面 $\mathbf{S}^{k}$, 则对 $n \geq 2$ 同样成立.
\item (\cite{klai33}、\cite{klai34}、\cite{klai19}. )若 $n = 3$, 则完整的 MKG 与 YM 系统在能量范数 $H^{2}(\mathbb{R}^{n})$ 下对大初值整体适定. 对临界幂非线性的 Yang--Mills--Higgs 系统结果亦成立.
\item (\cite{klai13}、\cite{klai21}. )若 $n = 3$, 则完整 MKG 系统在 $s > \frac{3}{4}$ 时局部适定, 在 $s > \frac{7}{8}$ 时整体适定.
\item (\cite{klai39}、\cite{klai47}. )若 $n = 4$, 则约化 MKG 与 YM 型系统在 $s > s_c = 1$ 时局部适定.\end{enumerate}
\end{theorem*}

\section*{2.1 主定理的动机}

考虑系统
\begin{equation}
\Box u=F(u,\partial u),\end{equation}
其中 $u : \mathbb{R}^{1+n} \to \mathbb{R}^N$, $\partial u = (\partial_t u, \partial_1 u, \ldots, \partial_n u)$, $F$ 为满足 $F(0) = 0$ 的光滑 $\mathbb{R}^N$ 值函数. 对这个方程有如下标准存在唯一性结果:
\setcounter{mtthm}{1}% 2.1 为节标题, 下一定理环境显示 2.2

\begin{theorem}\label{thm:2.2}
方程 (22) 对 $H^s \times H^{s-1}(\mathbb{R}^n)$ 中的初值局部适定, 只要 $s > \frac{n}{2} + 1$. 即对任意初值 $u[0] \in H^s(\mathbb{R}^n)$($s > \frac{n}{2} + 1$), 存在 $T > 0$ 与唯一解
\begin{equation}
u\in C^{0}([0,T],H^{s}(\mathbb{R}^{n}))\cap C^{1}([0,T],H^{s-1}(\mathbb{R}^{n})).\end{equation}
此外解连续(事实上是解析地)依赖于数据.\end{theorem}

\begin{remark}\label{rem:2.3}
这一非常经典结果的证明只基于简单的能量估计、嵌入 $L^{\infty} \subset H^{s}$ 的 Sobolev 估计, 以及一些简单的插值不等式. 而且证明可无大困难地推广到极一般的演化方程类, 包括拟线性波动方程组.\end{remark}

\begin{remark}\label{rem:2.4}
就初值的正则性假设而言, 该结果远非最优. 最优适定性的问题至关重要, 并在过去 10 年中产生了大量工作.\end{remark}

为在我们的例子(波动映射、Maxwell--Klein--Gordon 与 Yang--Mills 方程)背景下更好地理解最优局部适定性问题, 我们需要定义临界适定性(以下缩写为 WP)指数 $s_c$. 我们所有方程都有与之相伴的自然伸缩, 而 $s_c$ 是使初数据的 $\dot{H}^s \times \dot{H}^{s-1}$ 范数在此伸缩下不变的 $s$ 的唯一值. 例如, 若 $u$ 满足 (WM), 则对任意 $\lambda > 0$,
\[
u_{\lambda}(t,x)=u(\lambda t,\lambda x)
\]
也满足. 由于 $\|u_{\lambda}(t)\|_{H^{s}} = \lambda^{\frac{n}{2}-s}\|u(\lambda t)\|_{H^{s}}$, (WM) 的临界 WP 指数为 $s_{c} = \frac{n}{2}$.
同一原理适用于 (MKG)、(YM). 事实上它们临界 WP 指数均为 $s_{c} = \frac{n-2}{2}$.
据此我们给出如下猜想, 见 \cite{klai29}:
\begin{theorem*}[一般适定性猜想]
\begin{enumerate}[label=1.]
\item 对所有基本场论, 初值问题对 $H^{s}$($s > s_c$)中的初值局部适定.
\item 基本场论对所有具有小 $H^{s_c} \times H^{s_c-1}$ 范数的初数据弱$^{5}$ 整体适定.
\item 基本场论对 $H^{s} \times H^{s-1}$($s < s_c$)中的初值不适定.\end{enumerate}
\end{theorem*}

为证明 $s > s_c$ 时的局部适定性, 我们在适当的 Banach 空间中用 Picard 迭代. 考虑柯西问题
\begin{equation}
\Box u=\mathcal{N}(u),\qquad(u,\partial_{t}u)_{t=0}=(f,g).\end{equation}
第 $0$ 次迭代 $u_{0}$ 为解的无齐次部分:
\[
\Box u_{0}=0,\qquad(u,\partial_{t}u)_{t=0}=(f,g).
\]
后续迭代归纳地由
\[
u_{j+1}=u_{0}+\Box^{-1}\mathcal{N}(u_{j})\qquad (j \geq 0)
\]
给出, 其中 $\Box^{-1}$ 是把任一充分正则的 $F$ 映为方程 $\Box v = F$ 满足 $(v, \partial_t v)_{t=0} = 0$ 的解 $v$ 的算子.
若要通过迭代证明 $\Box u = \mathcal{N}(u)$ 在 $H^s \times H^{s-1}$ 中初值的局部解存在, 必须能证明迭代保持在数据空间中:
\[
f\in H^{s},\quad g\in H^{s-1}\implies u_{j}(t)\in H^{s},\quad\partial_{t}u_{j}(t)\in H^{s-1}
\]
对所有 $j \geq 0$ 与某区间 $(0,T)$ 中所有 $t$ 成立. 对 $j=0$, (23) 显然; 但 $j=1$ 的情形已提供有价值的见解.
\setcounter{mtthm}{4}% 2.3、2.4 为注, 下一定义环境显示 2.5

\begin{definition}\label{def:2.5}
若 (23) 对 $j = 1$ 与所有 $(f, g) \in H^s \times H^{s-1}$ 成立, 我们称第一迭代对 $H^s$ 中的初值 WP.\end{definition}

\begin{example}\label{ex:2.1}
考虑模型问题
\begin{equation}
\Box u=(\partial_{t}u)^{2},\end{equation}
其中 $u$ 为实值. 该方程与 (WM) 有相同的伸缩性质, 故 WP 指数为 $s_c = \frac{n}{2}$. 我们想求使第一迭代 $u_1$ 对 $H^s$ 中初值 WP 的 $s$ 的集合的下界. 一个涉及 Duhamel 原理的简单计算表明, 这归结为证明形如
\[
\int_{\mathbb{R}^{n}\times\mathbb{R}^{n}}K(\xi,\eta)f(\xi)g(\eta)h(\xi+\eta)\,d\xi d\eta\lesssim\|f\|_{L^{2}}\|g\|_{L^{2}}\|h\|_{L^{2}}
\]
的估计, 对所有 $f, g, h \in L^{2}(\mathbb{R}^{n})$ 成立, 其中
\[
K\left(\xi,\eta\right)=\frac{\langle\xi+\eta\rangle^{s-1}}{\langle\xi\rangle^{s-1}\langle\eta\rangle^{s-1}\left(1+\Delta_{\pm}(\xi,\eta)\right)},
\]
\[
\Delta_{+}=|\xi|+|\eta|-|\xi+\eta|,\qquad\Delta_{-}=|\xi+\eta|-|\xi|-|\eta|.
\]
这里记号 $\langle\cdot\rangle=1+|\cdot|$.
关于形如
\[
\int_{\mathbb{R}^{n}\times\mathbb{R}^{n}}\frac{f(\xi)g(\eta)h(\xi+\eta)}{\langle\xi\rangle^{a}\langle\eta\rangle^{b}\bigl(1+\Delta_{\pm}(\xi,\eta)\bigr)^{c}}\,d\xi d\eta\lesssim\|f\|_{L^{2}}\|g\|_{L^{2}}\|h\|_{L^{2}}
\]
的积分估计, 可得如下结果, 其中 $\Delta_{\pm}$ 由 (26) 给出.\end{example}

\begin{proposition}\label{prop:2.6}
设 $a, b, c \geq 0$. 则 (27) 成立当且仅当 $a + b + c > \frac{n}{2}$ 且 $c < \frac{n-1}{4}$.\end{proposition}

\begin{remark}\label{rem:2.7}
应当指出, 基于简单的集中型反例, 当 $a + b + c < \frac{n}{2}$ 或 $c \geq \frac{n-1}{4}$ 时估计失败.\end{remark}

把\propref{prop:2.6} 用于核 (25) 需要
\[
s-1+\min\left(1,\frac{n-1}{4}\right)>\frac{n}{2},\qquad \text{即}\ s>\max\left(\frac{n}{2},\frac{n+5}{4}\right).
\]
因此, 对 $n = 3$ 维的模型方程 $\Box u = (\partial_t u)^2$, 上例表明第一迭代对 $H^s$ 中初值在 $s > 2 = s_c + \frac{1}{2}$ 时 WP; 事实上可证当 $s \leq 2$ 时这不成立. 这应与 Lindblad \cite{klai49} 在 $n = 3$ 维的反例比较, 后者表明存在形如 $\Box u = q(\partial u)$ 的方程(其中 $q$ 是 $\mathbb{R}^4$ 上的二次型), 对 $H^2$ 中的数据不适定. 然而, 若二次型 $q$ 为零形式型, 则几乎可一路推进到临界 WP 指数 $s_c$. 下面两个例子在第一迭代层面验证了这一点.
\begin{example}\label{ex:2.2}
考虑方程$^{6}$
\begin{equation}
\Box u=Q_{0}(u,u),\end{equation}
其中 $u$ 为实值. 同样, 第一迭代的 WP 问题导至证明形如 (24) 的估计, 但由于算子 $Q_{0}$ 的特殊零结构, 奇异因子 $\Delta_{\pm}$ 从核的分母中完全消去. 事实上, $K$ 为
\[
K(\xi,\eta)=\langle\xi\rangle^{-s}+\langle\eta\rangle^{-s},
\]
故由\propref{prop:2.6}, 第一迭代对 $H^{s}$ 中数据 WP, 当 $s > s_{c} = \frac{n}{2}$.\end{example}

\begin{example}\label{ex:2.3}
考虑方程
\begin{equation}
\Box u=Q(u,u),\end{equation}
其中 $u$ 为向量值, $Q(u,u)$ 的第 $I$ 个分量是所有 $i,j,J,K$ 的 $Q_{ij}(u^{J},u^{K})$ 的线性组合. 如上例, 由于 $Q_{ij}$ 的零结构产生消去, 但此情形只消去 $\Delta_{\pm}$ 的半次幂. 事实上, 现在的 $K$ 为
\[
K(\xi,\eta)=\bigl(\langle\xi\rangle^{-s+\frac{1}{2}}+\langle\eta\rangle^{-s+\frac{1}{2}}\bigr)\left(1+\Delta_{\pm}(\xi,\eta)\right)^{-\frac{1}{2}},
\]
故第一迭代对 $H^{s}$ 中数据 WP, 当 $s > \max\left(\frac{n}{2}, \frac{n+3}{4}\right)$.
经显然的修改, 若考虑方程
\[
\Box u=Q(D^{-1}u,u),
\]
则发现第一迭代对 $H^{s}$ 中数据 WP, 当 $s > \max\left(\frac{n-2}{2}, \frac{n-1}{4}\right)$.\end{example}

高阶迭代 一旦第一迭代的分析完成, 真正困难的工作才开始. 为证明迭代 $u_{k+1}$ 满足 (23), 仅用所有先前 $u_j\,(j \leq k)$ 的 (23) 是不够的$^{7}$. 我们需要一个与第一讲所述的 Strichartz 及双线性估计密切相连、强大得多的函数空间.
记号  tempered 分布 $u \in \mathcal{S}'(\mathbb{R}^{n+1})$ 的傅里叶变换记为 $\mathcal{F}u$ 或 $\widehat{u}$, 维数任意. 在频率空间我们用坐标 $(\tau, \xi) = \Xi = (\Xi^0, \ldots, \Xi^n)$, 其中 $\tau \in \mathbb{R}$ 与 $\xi \in \mathbb{R}^n$ 分别对应于时间变量 $t$ 与空间变量 $x$. $\mathbb{R}^{1+n}$ 上的 Lorentz 内积为
\[
\langle\Xi;\widetilde{\Xi}\rangle=\sum_{\mu=0}^{n}\Xi_{\mu}\widetilde{\Xi}^{\mu}=-\Xi^{0}\widetilde{\Xi}^{0}+\sum_{j=1}^{n}\Xi^{j}\widetilde{\Xi}^{j},
\]
而波动算子 $\Box$ 的象征为 $-\langle \Xi;\Xi\rangle = \tau^2 - |\xi|^2$. 我们总以 $|\Xi|$ 表示欧氏范数 $|\Xi|^2 = \tau^2 + |\xi|^2$.
\begin{definition}\label{def:2.8}
设 $\Lambda^{\alpha}$、$\Lambda_{+}^{\alpha}$ 与 $\Lambda_{-}^{\alpha}$ 为如下乘子:
\begin{align}
\widehat{\Lambda^{\alpha}f}(\xi)&=\left(1+|\xi|^{2}\right)^{\alpha/2}\widehat{f}(\xi),\qquad
\widehat{\Lambda_{+}^{\alpha}u}(\Xi)=\left(1+|\Xi|^{2}\right)^{\alpha/2}\widehat{u}(\Xi),\nonumber\\
\widehat{\Lambda_{-}^{\alpha}u}(\Xi)&=\left(1+\frac{\langle\Xi;\Xi\rangle^{2}}{1+|\Xi|^{2}}\right)^{\alpha/2}\widehat{u}(\Xi).\end{align}
\end{definition}

我们定义空间 $H^{s,\theta}(\mathbb{R}^{n+1})$, 它适配于 $\mathbb{R}^{1+n}$ 上的波动算子, 正如 $H^{s}(\mathbb{R}^{n})$ 适配于 $\mathbb{R}^{n}$ 上的 Laplacian.
\begin{definition}\label{def:2.9}
对 $s, \theta \in \mathbb{R}$, 定义
\[
H^{s,\theta}=\left\{u\in\mathcal{S}^{\prime}(\mathbb{R}^{n+1}):\Lambda^{s}\Lambda_{-}^{\theta}u\in L^{2}(\mathbb{R}^{n+1})\right\},
\]
其范数为 $\|u\|_{s,\theta}=\|\Lambda^{s}\Lambda_{-}^{\theta}u\|_{L^{2}}$.\end{definition}

$H^{s,\theta}$ 与具 $H^s$ 中数据的齐次波动方程解空间之间有极简单的联系. 事实上, 每个 $u \in H^{s,\theta}$ 都可写成形式
\[
u(t)=\frac{1}{2\pi}\int_{-\infty}^{\infty}\frac{e^{i t\lambda}u_{\lambda}(t)}{(1+\left|\lambda\right|)^{\theta}}\,d\lambda\qquad(H^{s}\text{ 值积分}),
\]
其中 $\{u_\lambda\}_{\lambda\in\mathbb{R}}$ 是 (1) 的一族单参数解, 数据属于 $H^s$; 即 $\Box u_\lambda = 0$ 且 $(u_\lambda, \partial_t u_\lambda)_{t=0} = (f_\lambda, 0)$, 其中 $\lambda \mapsto f_\lambda$ 属于 $L^2(\mathbb{R}, H^s)$. 而且 $\|u\|_{s,\theta}^2 = \int \|f_\lambda\|^2 d\lambda$.
(28) 的一个重要推论如下:
\begin{theorem*}[原理]
对具 $H^{s}$ 中数据的齐次波动方程解的一个线性或多次线性时空估计, 蕴含对 $\theta > \frac{1}{2}$ 时 $H^{s,\theta}$ 中元素的相应估计.\end{theorem*}

空间 $H^{s,\theta}$ 提供了证明主定理的粗略框架. 事实上它们只适用于证明主定理的 i) 部分. 所有进一步进展都要求对这些空间作实质性精化. 关于 $H^{s,\theta}$ 空间及其修正, 连同第一讲的双线性与 Strichartz 估计与上述原理, 如何用于证明主定理中 $s > s_c$ 的适定性结果, 我们请读者参阅我们的综述 \cite{klai45}. 情形 $s = s_c$ 在 T. Tao 的近期工作之前完全开放, 这将在下一节描述.
\section*{2.10 临界正则性的波动映射}

下面我将讨论 T. Tao 最近得到的关于临界正则性波动映射的突破性结果. 他的结果有力地展示了物理空间中几何思想与傅里叶分析技巧(如仿微分法、Strichartz 与双线性估计)之间的相互作用. 因此它非常契合本讲稿的框架.
他的结果适用于波动映射 $\phi : \mathbb{R}^{n+1} \to S^N$, $S^N$ 为欧氏空间 $\mathbb{R}^{N+1}$ 中的单位球面, 对任意维数 $n \geq 2$. 维数 $n \geq 5$ 的情形容易得多, 从而可专注于 Tao 的主要新思想------规范重整化(gauge renormalization). 事实上我将介绍 Tao 的 $n \geq 5$ 结果由 Rodnianski 与我和 \cite{klai41} 推广到一般目标流形类的一个推广.
设 $\phi: \mathbb{R}^{n+1} \longrightarrow (\mathcal{N}, h)$, 其中 $(\mathcal{N}, h)$ 为维数 $N$ 的黎曼流形. 这里 $\mathbb{R}^{n+1}$ 表示赋予度量 $m = \mathrm{diag}(-1, 1, \ldots, 1)$ 的标准 Minkowski 空间. 记 $\nabla$ 为 $\mathcal{N}$ 的切丛 $T\mathcal{N}$ 上的 Levi-Civita 联络. 在 $\mathcal{N}$ 上的局部坐标 $y^I\,(I = 1, \ldots, N)$ 中, 波动映射方程取熟知形式
\[
\partial^{\alpha}\partial_{\alpha}\phi^{I}+\Gamma_{J K}^{I}\partial_{\beta}\phi^{J}\partial_{\gamma}\phi^{K}m^{\beta\gamma}=0,
\]
其中 $\Gamma_{JK}^{I}$ 为 $\mathcal{N}$ 上 Levi-Civita 联络 $\nabla$ 的 Christoffel 系数, $\phi^{I}\,(I = 1, \ldots, N)$ 为映射 $\phi$ 在 $\mathcal{N}$ 局部坐标中的分量. 令 $e_{a} = e_{a}^{I} \frac{\partial}{\partial y^{I}}$ 为一组正交标架向量场, $\omega^{a} = \omega_{I}^{a} dy^{I}$ 为对应的对偶 1-形式基, $\omega^{a}(e_{b}) = \delta_{b}^{a}$. 由于 $h(e_{a}, e_{b}) = \delta_{ab}$, 推得 $h_{IJ} = \sum_{a} \omega_{I}^{a} \omega_{J}^{a}$. 定义
\[
\phi_{\alpha}^{a}=\omega_{I}^{a}\partial_{\alpha}\phi^{I},
\]
其中 $\phi^{I}$ 为映射 $\phi$ 相对于 $\mathcal{N}$ 上局部坐标 $y^{I}$ 的分量. 显然 $\partial_{\alpha}\phi^{I} = e_{a}^{I}\phi_{\alpha}^{a}$. 对 $\mathcal{N}$ 上函数 $F$ 记 $\partial_{\alpha}F(\phi) = \frac{\partial}{\partial y^{I}}F(\phi)\partial_{\alpha}\phi^{I} = \frac{\partial}{\partial y^{I}}F(\phi)e_{d}^{I}\phi_{\alpha}^{d} = F_{,d}(\phi)\phi_{\alpha}^{d}$, 其中 $F_{,d} = e_{d}(F)$.
我们易验证与波动映射 $\phi$ 相伴的函数 $\phi_{\alpha}^a = \langle \partial_\alpha \phi, e_a \rangle$ 满足$^{8}$ 如下散度--旋度系统:
\[
\partial_{\beta}\phi_{\alpha}^{a}-\partial_{\alpha}\phi_{\beta}^{a}=C_{b c}^{a}\phi_{\alpha}^{b}\phi_{\beta}^{c},
\]
\[
\partial^{\alpha}\phi_{\alpha}^{a}=-\Gamma_{b c}^{a}\phi_{\beta}^{b}\phi_{\gamma}^{c}m^{\beta\gamma},
\]
其中 $C_{bc}^{a}$ 与 $\Gamma_{bc}^{a}$ 分别为标架的结构系数与联络系数,
\[
[e_{b},e_{c}]=C_{b c}^{a}e_{a},
\]
\[
\nabla_{e_{b}}e_{c}=\Gamma_{b c}^{a}e_{a}.
\]
由公式 $[e_b, e_c] = \nabla_{e_b} e_c - \nabla_{e_c} e_b$ 推得
\[
C_{b c}^{a}=\Gamma_{b c}^{a}-\Gamma_{c b}^{a}.
\]
由于标架 $e_a$ 正交, $\Gamma^a_{bc} = \langle \nabla_{e_b} e_c ; e_a \rangle = - \langle e_c ; \nabla_{e_b} e_a \rangle$, 故
\[
\Gamma_{b c}^{a}=-\Gamma_{b a}^{c}.
\]
此外,
\[
\Gamma_{b c}^{a}=\frac{1}{2}\bigl(C_{b c}^{a}+C_{a c}^{b}+C_{a b}^{c}\bigr).
\]

\setcounter{mtthm}{10}% 2.10 为节标题, 下一定义环境显示 2.11

\begin{definition}\label{def:2.11}
若黎曼流形 $\mathcal{N}$ 上存在一组正交标架 $(e_a)_{a}$ 使得结构系数 $C_{bc}^a$ 及其标架导数 $C_{bc,d_1d_2...d_k}^a$ 在 $\mathcal{N}$ 上一致有界, 则称 $\mathcal{N}$ 具有"有界可平行化"(bounded parallelizable)结构.\end{definition}

\begin{remark}\label{rem:2.12}
有界可平行化流形有大量例子. 首先, 在任何李群上我们可构造一组左不变向量场的正交基 $e_a$, 相对它结构常数 $C_{bc}^a$ 为常数$^{9}$. 常负曲率流形 $\mathbf{H}^N\,(N > 2)$ 是有界可平行化的. 而且 $\mathbf{H}^2$(即双曲平面)是一个李群, 见 \cite{klai11} 第 3.1 节的讨论. 此外任何紧黎曼流形都可等距嵌入为一个有界可平行化黎曼流形的全测地子流形, 见 \cite{klai11}.\end{remark}

\begin{proposition}\label{prop:2.13}
设 $\mathcal{N}$ 为黎曼流形, $\phi: \mathbb{R}^{n+1} \to \mathcal{N}$ 为波动映射. 1-形式 $\phi_\alpha^a = \langle \partial_\alpha \phi, e_a \rangle$ 满足方程 (31)、(32) 以及如下波动方程组:
\[
\Box\Phi=-2R_{\mu}\cdot\partial^{\mu}\Phi+E,
\]
其中 $\Phi = (\phi_{\alpha}^{a})$, $R_{\mu} = (R_{b\mu}^{a})_{a,b=1}^{N}$, $R_{b\mu}^{a} = \Gamma_{cb}^{a}\phi_{\mu}^{c}$. 分量 $E = (E_{\alpha}^{a})$ 是 $\Phi = (\phi_{\alpha}^{a})$ 各分量的三次齐次多项式, 其系数只依赖于结构函数 $C_{bc}^{a}$ 及其关于标架的导数 $C_{bc,d}^{a}$.\end{proposition}

\begin{remark}\label{rem:2.14}
必须指出矩阵 $R_{\mu}$ 是反对称的, 即
\[
R_{b\mu}^{a}=-R_{a\mu}^{b}.
\]
这是 (??) 的直接推论. 这表明著名的"Helein 技巧"------在 $N$ 为标准球面的特例下将波动映射方程反对称化, 这一技巧在 Tao 关于标准球面的工作 \cite{klai70}、\cite{klai71} 中起根本作用------事实上源于任何黎曼流形上相对正交标架的联络系数的一般特征.\end{remark}

我们研究波动映射满足初值问题
\[
\phi(0)=\varphi,\qquad\partial_{t}\phi(0)=\psi=\psi_{0}^{a}e_{a}
\]
的演化, 其中 $\varphi$ 为从 $\mathbb{R}^n$ 取值于 $\mathcal{N}$ 的任意光滑映射, $\psi = \psi_0^a e_a$ 为从 $\mathbb{R}^{n+1}$ 到 $T\mathcal{N}$ 的任意光滑映射. 令 $\varphi_i^a = \langle \partial_i \varphi, e_a \rangle$.
\begin{definition}\label{def:2.15}
若所有分量 $\varphi_i^a$、$\psi_i^a$ 属于空间 $\dot{H}^{s-1}(\mathbb{R}^n)$(相应地 $H^{s-1}(\mathbb{R}^n)$), 则称初数据 $\phi[0] = (\varphi, \psi)$ 属于 Sobolev 空间 $\dot{H}^s(\mathbb{R}^n)$(相应地 $H^s(\mathbb{R}^n)$). 我们记
\[
\|\phi[0]\|_{\dot{H}^{s}}=\sum_{a,i}\left(\|\varphi_{i}^{a}\|_{\dot{H}^{s-1}}+\|\psi_{i}^{a}\|_{\dot{H}^{s-1}}\right),
\]
对 $\|\phi[0]\|_{H^{s}}$ 类似.\end{definition}

\begin{theorem}\label{thm:2.16}
设 $\mathcal{N}$ 为装备有界可平行化结构的黎曼流形. 设 $n \geq 5$ 且初数据 $\phi[0] = (\varphi, \psi = \psi_i^a e_a)$ 对某 $\frac{n}{2} < s$ 属于 $H^s$. 我们还作临界小性假设:
\[
\|\phi[0]\|_{\dot{H}^{\frac{n}{2}}}\leq\varepsilon.
\]
则具初数据 $\phi[0]$ 的波动映射 $\phi$ 可在 $H^{s}$ 范数下唯一地整体延拓到时刻 $t$.\end{theorem}

\begin{remark}\label{rem:2.17}
该定理把 \cite{klai70} 的结果从目标流形为标准球面推广到有界可平行化流形. 对维数的限制 $n \geq 5$ 与 \cite{klai70} 中相同; 这使我们只需依赖 Strichartz 估计. 对标准球面的情形, 维数限制由 \cite{klai71} 借助双线性估计$^{10}$、$H^{s,\theta}$ 空间以及 \cite{klai73} 的精化方法移除. 即便有 \cite{klai71}, 将我们的结果推广到二维看来也并非易事$^{11}$.\end{remark}

主定理的证明依赖于 $H^{s}$($s > \frac{n}{2}$)中的局部适定性结果. 我们陈述如下精确结果:
\begin{theorem}\label{thm:2.18}
设初数据 $\phi[0] \in H^s(\mathbb{R}^n)$ 对某 $s \geq s_0 > \frac{n}{2}$. 存在 $T > 0$(只依赖于 $\|\phi[0]\|_{H^{s_0}}$ 的大小)与系统 (31)、(32) 在板 $[0, T] \times \mathbb{R}^n$ 上的唯一解 $\phi$, 满足
\[
\left\|\phi[t]\right\|_{H^{s}}\leq C\|\phi[0]\|_{H^{s}}
\]
对所有 $t \in [0, T]$ 成立, 其中常数 $C$ 只依赖于 $T$、$s$、$s_0 - \frac{n}{2}$ 与 $n$.\end{theorem}

\begin{proof}
我们用 \cite{klai70} 的 Littlewood--Paley 记号. 对函数 $\phi(t, x)$ 记投影
\[
P_k\phi(t, x) = \int e^{ix \cdot \xi} \chi(2^{-k}\xi)\hat\phi(t, \xi) \, d\xi,
\]
其中 $\hat\phi(t, \xi)$ 为 $\phi$ 的空间傅里叶变换, $\chi(\xi) = \eta(\xi) - \eta(2\xi)$, $\eta$ 为支在 $|\xi| \leq 2$ 上、在球 $|\xi| \leq 1$ 上等于 1 的非负光滑鼓包函数. 故 $\chi(\xi)$ 支在 $\{\frac{1}{2} \leq |\xi| \leq 2\}$, 且对 $\xi \neq 0$ 有 $\sum_{k \in \mathbf{Z}} \chi(2^{-k}\xi) = 1$. 又定义 $P_{\leq k} = \sum_{l \leq k} P_l$, 以及对任意区间 $I \subset \mathbf{Z}$ 显然地定义 $P_I$.
记号: 我们将频繁使用记号 $A \lesssim B$ 表示对某不依赖于估计中任何重要参数的常数 $c > 0$ 有 $A \leq cB$.
依 \cite{klai70} 引入记号
\[
\|\Phi\|_{S_{k}}=\sup_{q,r\in\mathcal{A}}2^{k\left(\frac{1}{q}+\frac{n}{r}-1\right)}\left(\|\Phi\|_{L_{t}^{q}L_{x}^{r}}+2^{-k}\|\partial_{t}\Phi\|_{L_{t}^{q}L_{x}^{r}}\right),
\]
其中 $\mathcal{A} = \{(q, r): 2 \leq q, r \leq \infty,\ \frac{1}{q} + \frac{n-1}{2r} \leq \frac{n-1}{4}\}$ 为容许 Strichartz 指数的集合. 回顾:\end{proof}

\begin{proposition}[Strichartz 型估计]\label{prop:2.19}
对任意固定整数 $k$ 与 $\mathbb{R} \times \mathbb{R}^n$ 上函数 $\phi(t, x)$, 若其 $\hat{\phi}(t, \xi)$ 的支包含于二进区域 $2^{k-1} \leq |\xi| \leq 2^{k+1}$, 则有估计
\[
\|\phi\|_{S_{k}}\lesssim\|\phi(0)\|_{H^{\frac{n-2}{2}}}+\|\partial_{t}\phi(0)\|_{H^{\frac{n-4}{2}}}+2^{k\frac{n-4}{2}}\|\Box\phi\|_{L_{t}^{1}L_{x}^{2}}.
\]
\end{proposition}

下面我们回顾 \cite{klai70} 中给出的频率包络定义.
\begin{definition}\label{def:2.20}
频率包络是一个 $l^2$ 序列 $c = (c_k)_{k \in \mathbf{Z}}$ 满足
\[
c_{k}\lesssim2^{\sigma|k-k^{\prime}|}c_{k^{\prime}}
\]
对所有 $k, k' \in \mathbf{Z}$ 成立. 这里 $\sigma$ 为固定的正常数; 如 \cite{klai70} 取 $0 < \sigma < \frac{1}{2}$. 此外我们需要 $0 < \sigma < \frac{n-4}{4}$ 与 $0 < \sigma < \frac{n-3}{4(n-1)}$.\end{definition}

若 $\mathbb{R}^n$ 上函数 $f$ 的 $\hat{H}^s$ 范数位于包络 $c$ 之下, 是指对所有 $k \in \mathbf{Z}$ 有 $\|P_k f\|_{\dot{H}^s} \leq c_k$. 记为 $f \ll_s c$, 无歧义时简记 $f \ll c$. 易验证若 $\|f\|_{\dot{H}^s} \leq \varepsilon$, 则存在包络 $c \in l^2$ 使 $\|c\|_{l^2} \lesssim \varepsilon$ 且 $f \ll_s c$. 事实上我们可取 $c_k = \sum_{k' \in \mathbf{Z}} 2^{-\sigma |k - k'|} \|P_k\|_{\dot{H}^s}$.
\begin{definition}\label{def:2.21}
固定 $0 < \sigma < \min(\frac{1}{2}, \frac{n-4}{4}, \frac{n-3}{4(n-1)})$ 与频率包络 $c$. 若相对我们的标架 $e_a$, 对所有 $k \in \mathbf{Z}$ 有
\[
\left\|P_{k}\phi[0]\right\|_{\dot{H}^{\frac{n}{2}}}\leq c_{k},
\]
则称初数据 $\phi[0] = \left( \phi(0) = \varphi, \partial_t \phi(0) = \psi = \psi_0^a e_a \right)$ 位于 $c$ 之下. 简记为 $\phi[0] \ll c$.\end{definition}

我们的主定理的证明可轻易化为如下:
\begin{proposition}[主命题]\label{prop:2.22}
设 $c$ 为频率包络$^{12}$, 满足 $\|c\|_{l^2} \leq \varepsilon$, $0 < T < \infty$, 且 $\Phi = (\phi_\alpha^a)$ 满足方程 (31)、(32), 从而也满足 (34). 依\defref{def:2.21}, 设初数据满足小性条件 $\phi[0] \ll c$. 再设自举(bootstrap)假设
\[
\left\|P_{k}\Phi\right\|_{S_{k}([0,T]\times\mathbb{R}^{n})}\leq2C c_{k}
\]
对所有 $k \in \mathbf{Z}$ 成立. 则事实上对充分小的 $\varepsilon$ 与所有 $k \in \mathbf{Z}$,
\[
\|P_{k}\Phi\|_{S_{k}([0,T]\times\mathbb{R}^{n})}\leq C c_{k}.
\]
\end{proposition}

\begin{remark}\label{rem:2.23}
由于我们的方程与小性条件 $\phi[0] \ll c$ 的尺度不变性, 只需对 $k = 0$ 证明 (40). 令 $\Psi = P_0\Phi$. 我们需证明
\[
\|\Psi\|_{S_{0}([0,T]\times\mathbb{R}^{n})}\leq C c_{0}.
\]
\end{remark}

\begin{remark}\label{rem:2.24}
为证明 (41), 我们想把\propref{prop:2.19} 应用于对 (34) 作投影 $P_{0}$ 所得的方程, 即
\[
\Box\Psi=P_{0}(R_{\mu}\cdot\partial^{\mu}\Phi+E).
\]
然而 Strichartz 不等式的直接应用不会奏效. 事实上依\propref{prop:2.19},
\[
\|\Psi\|_{S_{0}}\leq c_{0}+\|P_{0}(R_{\mu}\cdot\partial^{\mu}\Phi+E)\|_{L_{t}^{1}L_{x}^{2}}.
\]
三次项 $E$ 不成问题, 困难出现在我们试图更精细地估计 $P_0(R_\mu \cdot \partial^\mu \Phi)$------更精确地, 出现在 $R$ 的低频与同 $\Psi$ 频率相当的 $\Phi$ 频率相互作用的部分. 最危险的项形如 $\tilde{R} \cdot \partial \Psi$, 其中 $\tilde{R} = P_{\leq -10} R$. 为相对于可用的 Strichartz 范数估计 $\|\tilde{R} \cdot \partial \Psi\|_{L_t^1 L_x^2}$, 我们被迫取 $\Psi$ 在能量范数 $L_t^\infty L_x^2$ 中. 这留下 $\tilde{R}$ 在范数 $L_t^1 L_x^\infty$ 中, 而对此我们没有 Strichartz 估计.\end{remark}

\begin{definition}\label{def:2.25}
下面我们用 Tao 的约定, 称 $[0, T] \times \mathbb{R}^n$ 上任一函数或矩阵值函数 $F$ 为可接受误差(acceptable error), 如果
\[
\|F\|_{L_{t}^{1}L_{x}^{2}([0,T]\times\mathbb{R}^{n})}\leq C^{3}\varepsilon c_{0}.
\]
\end{definition}

\begin{proposition}\label{prop:2.26}
矩阵值函数 $P_{0}\Phi = \Psi$ 满足方程
\[
\Box\Psi=-2\tilde{R}_{\mu}\cdot\partial^{\mu}\Psi+\text{error},
\]
其中 $\tilde{R}_{\mu} = P_{\leq -10} R_{\mu} = \Gamma \cdot \tilde{\Phi}_{\mu}$, $\tilde{\Phi}_{\alpha} = P_{\leq -10} \Phi_{\alpha}$. 这里"error"指 (42) 意义下的可接受误差项.\end{proposition}

\begin{proposition}\label{prop:2.27}
矩阵值函数 $\Psi$ 满足形如
\[
\Box\Psi=-2\partial_{\mu}\tilde{\Delta}\cdot\partial^{\mu}\Psi+\text{error}
\]
的方程, 其中势 $\Delta$ 满足如下性质:
\begin{enumerate}[label=\roman*.)]
\item $N \times N$ 矩阵 $\tilde{\Delta}$ 反对称, 即 $\tilde{\Delta}^t = -\tilde{\Delta}$. $\tilde{\Delta}$ 各分量的空间傅里叶变换支在 $|\xi| \leq 2^{-10}$.
\item 对任意 $\tilde{\Delta}_{k} = P_{k}\tilde{\Delta}$ 成立估计:
\[
\|\tilde{\Delta}_{k}\|_{S_{k}}\lesssim 2^{-k}C c_{k},\qquad
\|\partial\tilde{\Delta}_{k}\|_{S_{k}}\lesssim C c_{k},\qquad
\|\Box\tilde{\Delta}_{k}\|_{S_{k}}\lesssim2^{k}C c_{k}.
\]
\item 令 $\bar{R}_{\mu} = \tilde{R}_{\mu} - \partial_{\mu} \tilde{\Delta}$. 对所有 $P_{k} \bar{R}$ 成立:
\[
\|P_{k}\bar{R}\|_{L_{t}^{1}L_{x}^{\infty}}\lesssim C^{2}c_{k}^{2},\qquad
\|P_{k}\bar{R}\|_{L_{t}^{\infty}L_{x}^{\infty}}\lesssim 2^{k}C^{2}c_{k}^{2}.
\]
\end{enumerate}
\end{proposition}

设 $M$ 为依赖于 $T$ 的大整数, 将在下面选取. 定义实 $N \times N$ 矩阵值函数 $U$ 为
\[
U=I+\sum_{-M<k\leq-10}U_{k},
\]
其中 $U_{k}$ 归纳定义如下:
\[
U_{k}=0\quad\text{对所有}\ k<-M;\qquad U_{-M}=I;\qquad
U_{k}=\tilde{\Delta}_{k}\cdot U_{<k}\quad\text{对所有}\ -M<k\leq-10,
\]
其中 $U_{<k} = \sum_{k' < k} U_{k'}$. 由于矩阵 $\tilde{\Delta}_k = P_k\tilde{\Delta}$ 反对称, 得到恒等式
\[
U_{k}^{t}\cdot U_{<k}+U_{<k}^{t}\cdot U_{k}=0,
\]
从而
\[
U_{<k}^{t}\cdot U_{<k}-I=\sum_{k^{\prime}<k}U_{k^{\prime}}^{t}\cdot U_{k^{\prime}}.
\]
利用此恒等式可归纳证明
\[
\|U_{<k}\|_{L_{t}^{\infty}L_{x}^{\infty}}\leq 2,\qquad
\|U_{k}\|_{L_{t}^{\infty}L_{x}^{\infty}}\lesssim C c_{k}\quad (k>-M),
\]
以及
\[
\|U_{k}\|_{L_{t}^{2}L_{x}^{\infty}}\lesssim C2^{-k/2}c_{k}\quad (k>-M).
\]
此外,
\[
\|\partial U_{<k}\|_{L_{t}^{\infty}L_{x}^{\infty}}\lesssim 2^{k}C^{2}c_{k},\qquad
\|\partial U_{k}\|_{L_{t}^{\infty}L_{x}^{\infty}}\lesssim 2^{k}C^{2}c_{k},
\]
以及
\[
\|\partial U_{<k}\|_{L^{2}_{t} L^\infty_{x}}\lesssim 2^{k/2}C c_{k},\qquad
\|\partial U_{k}\|_{L^{2}_{t} L^\infty_{x}}\lesssim 2^{k/2}C c_{k}.
\]

\begin{proposition}\label{prop:2.28}
设 $\varepsilon$ 充分小(依赖于 $C$), 且 $M$ 充分大(依赖于 $T, C, \varepsilon$). 则矩阵 $U$ 满足如下性质:
\begin{enumerate}[label=\roman*.)]
\item 近似正交性:
\[
\|U^{t}U-I\|_{L_{t}^{\infty}L_{x}^{\infty}},\quad\|\partial(U^{t}U-I)\|_{L_{t}^{\infty}L_{x}^{\infty}}\lesssim C^{2}\varepsilon.
\]
特别地, 对小 $\varepsilon$, $U$ 可逆且
\[
\|U\|_{L_{t}^{\infty}L_{x}^{\infty}},\quad\|U^{-1}\|_{L_{t}^{\infty}L_{x}^{\infty}}\lesssim1.
\]
\item 近似规范条件:
\[
\|\partial_{\mu}U-\partial_{\mu}\tilde{\Delta}\cdot U\|_{L_{t}^{1}L_{x}^{\infty}}\lesssim C^{2}\varepsilon.
\]
\item 此外,
\[
\|\partial U\|_{L_{t}^{\infty}L_{x}^{\infty}},\quad\|\partial U\|_{L_{t}^{1}L_{x}^{\infty}}\lesssim C^{2}\varepsilon,\qquad
\|\Box U\|_{L_{t}^{2}L_{x}^{n-1}}\lesssim C^{2}\varepsilon.
\]
\end{enumerate}
\end{proposition}

依 \cite{klai70} 我们现在可以进行规范变换
\[
\Psi=U\cdot W.
\]
$W$ 满足方程
\[
\Box W=-2U^{-1}(\partial_{\mu}U-\partial_{\mu}\tilde{\Delta}\cdot U)\partial^{\mu}W
-2U^{-1}\partial_{\mu}\tilde{\Delta}\cdot(\partial^{\mu}U)U^{-1}\Psi-U^{-1}(\Box U)U^{-1}\Psi+\text{error}.
\]
由\propref{prop:2.28} 推得:
\begin{proposition}\label{prop:2.29}
矩阵值函数 $W$ 满足形如
\[
\Box W=\text{error}
\]
的方程. 因此, 若 $\varepsilon$ 充分小,
\[
\|\Psi\|_{S_{0}}\lesssim\|W\|_{S_{0}}\leq\|\Psi[0]\|_{H^{\frac{n-2}{2}}}+C C^{3}\varepsilon c_{0}\leq C c_{0}.
\]
这正是 (41), 从而结束主命题的证明.\end{proposition}

\chapter{第三讲\quad 广义相对论与拟线性波动方程}

我们从研究拟线性波动方程最优适定性问题的主要动机开始.
\section*{3.1 广义相对论中的柯西问题}

爱因斯坦广义相对论的主要对象是时空. 它可定义为可微、定向的四维 Lorentz 流形 $(\mathcal{M}, g)$ 的等价类. 称两个 Lorentz 流形 $(\mathcal{M}, g)$、$(\mathcal{M}', g')$ 等价, 若存在微分同胚 $\mathcal{M} \longrightarrow \mathcal{M}'$ 使 $g' = \Phi_* g$. 时空就是此类 Lorentz 流形的一个等价类. 时空度量 $g$ 须满足爱因斯坦场方程
\[
R_{\alpha\beta}-\frac{1}{2}g_{\alpha\beta}R=T_{\alpha\beta},
\]
其中 $R_{\alpha\beta}$ 为 Ricci 曲率, $R$ 为度量的标量曲率, $T_{\alpha\beta}$ 为定义在 $(\mathcal{M}, g)$ 上某物质场的能量--动量张量. 我们限于真空特例, 即 $T\equiv0$, 此时方程变为
\[
R_{\alpha\beta}=0.
\]

爱因斯坦方程的演化特征比其他方程微妙得多, 因为并不存在时间变量 $t$ 的内在定义. 因此谈论初值问题似乎同样含混. 不过它可如下定义:
\setcounter{mtthm}{1}
\begin{definition}\label{def:3.2}
初数据集是一个三元组 $(\Sigma, \bar{g}, \bar{k})$, 由一个三维完备黎曼流形 $(\Sigma, \bar{g})$ 与一个在 $\Sigma$ 上的 2-协变对称张量 $k$ 组成, 满足约束方程:
\[
\bar{\nabla}^{j}\bar{k}_{i j}-\bar{\nabla}_{i}\mathrm{tr}\,\bar{k}=0,\qquad
\bar{R}-|\bar{k}|^{2}+(\mathrm{tr}\,\bar{k})^{2}=0,
\]
其中 $\bar{\nabla}$ 为协变导数, $\bar{R}$ 为 $(\Sigma, \bar{g})$ 的标量曲率. 标量 $\bar{k}$ 为 $\bar{k}$ 相对 $\bar{g}$ 的迹.\end{definition}

\setcounter{mtthm}{2}
\begin{definition}\label{def:3.3}
若初数据集对应于 Minkowski 空间中具诱导度量与第 II 基本形式的完备类空超曲面, 则称其为平坦(或平凡)的. 若在某坐标系 $(x^1, x^2, x^3)$ 定义于 $\Sigma$ 上无穷远邻域$^{13}$ 中, 使得度量 $\bar{g}$ 趋于欧氏度量、$\bar{k}$ 趋于零$^{14}$, 则称初数据集为渐近平坦(AF):
\[
r=\sqrt{(x^{1})^{2}+(x^{2})^{2}+(x_{3})^{2}}\to\infty.
\]
\end{definition}

\setcounter{mtthm}{3}
\begin{definition}\label{def:3.4}
初数据集 $(\Sigma, \bar{g}, \bar{k})$ 的柯西发展是一个时空流形 $(\mathcal{M}, g)$, 满足爱因斯坦方程以及一个嵌入$^{15}$ $i: \Sigma \longrightarrow \mathcal{M}$, 使得 $\bar{g}, \bar{k}$ 为 $\Sigma$ 在 $\mathcal{M}$ 中的诱导第 I、第 II 基本形式. 为保证解对数据的唯一依赖, 还要求该发展是整体双曲的$^{16}$.\end{definition}

爱因斯坦方程演化特征的微妙性, 使得对初值问题的局部存在唯一性(对很大一类演化方程已满意解决)成为首要问题. 对爱因斯坦方程, 此类结果最早由 Y. Choquet-Bruhat \cite{klai7} 借助波坐标得到, 波坐标使她能把爱因斯坦方程写成非线性波动方程组的形式, 从而可应用只基于能量估计与 $L^\infty \subset H^s$ Sobolev 不等式的标准局部存在理论.
\setcounter{mtthm}{4}
\begin{definition}\label{def:3.5}
若坐标 $x^\alpha$ 满足 $\Box_g x^\alpha = 0$, 则称其为波坐标. 给定这样的坐标系, 爱因斯坦真空方程取形式
\[
g^{\alpha\beta}\partial_{\alpha}\partial_{\beta}g_{\mu\nu}=N_{\mu\nu}(g_{\alpha\beta},\partial g_{\alpha\beta}).
\]
\end{definition}

对非线性微分方程组, 局部存在唯一性结果通过简单的延拓论证导出整体结果. 对爱因斯坦方程表述同类结果稍显微妙; 这由 Y. Choquet-Bruhat 与 R.P. Geroch 在 \cite{klai8} 中完成. 该结果断言: 对每个初数据集存在唯一的极大未来发展. 因此, 由演化方法从特定初数据集出发构造的任何发展, 必包含于其极大发展中. 这可以说解决了广义相对论中整体$^{17}$ 存在唯一性问题; 所有进一步问题可说关乎极大发展的定性性质. 核心问题是奇点的存在性与特征. 首先, 我们可定义正则极大发展为完备者, 即所有未来类时与零测地线都完备. Penrose 的一个著名定理断言, 在初数据集的某些相当极端的条件下$^{18}$, 相应的未来极大发展必不完备. 此外, 除平坦 Minkowski 空间本身外, 爱因斯坦真空方程的所有显式 AF 解(如 Kerr 族)都有奇点. 与 Penrose 定理相对的是如下完备性结果, 见 \cite{klai12}.
\setcounter{mtthm}{5}
\begin{theorem}[Minkowski 空间的全局稳定性]\label{thm:3.6}
任何充分接近平凡数据的渐近平坦初数据集, 具有完备的极大未来发展 $(M,g)$. 而且时空 $M$ 在如下意义下接近 Minkowski 空间: 其曲率张量沿任何类时或零测地线趋于零. 然而其因果结构渐近非平凡.\end{theorem}

\begin{problem*}[闵可夫斯基空间的强稳定性问题]
是否存在一个尺度不变的小性条件$^{20}$, 使得所有初数据集 $(\Sigma, g, k)$ 满足它的发展都有完备的极大未来发展? 
\end{problem*}

\setcounter{mtthm}{6}
\begin{remark}\label{rem:3.7}
\thmref{thm:3.6} 的旧稳定性结果与强稳定性之间的关系, 可与 Tao 的结果 \cite{klai71} 之后波动映射的情形相比. 事实上 Sideris \cite{klai63} 的一个较早结果对含大量导数的加权 Sobolev 范数中的小数据给出整体存在. Tao 的新结果只要求含初值映射最少导数的尺度不变范数的小性.\end{remark}

\setcounter{mtthm}{7}
\begin{remark}\label{rem:3.8}
Minkowski 时空的强稳定性本身应被视为通向证明著名的 Penrose 宇宙监督猜想的更高目标的中间步骤. 我们在第三讲的补充中陈述这些猜想.\end{remark}

证明 Minkowski 空间的强稳定性是一个长期目标. 下面我们提出一个远为简单、但仍极具挑战性的中间猜想, 它启发下一节的主题.
\begin{theorem*}[有限 $L^{2}$ 曲率猜想]
Bruhat--Geroch 结果可推广到初数据集 $(\Sigma, g, k)$ 满足
\[
R(g)\in L^{2}_{(loc)}(\Sigma)\qquad \text{且}\qquad k\in H^{1}_{(loc)}(\Sigma).
\]
\end{theorem*}

$L^2$ 有界性猜想是我们此刻就 Minkowski 空间强稳定性猜想所能期望的最合理目标. 它可与 $\mathbb{R}^{1+3}$ 中 YM 方程局部(在此情形也是整体的! )适定性的证明相比, 见 \cite{klai34}. 但它要复杂得多. 我们或许先满足于一个更简单的目标: 在 $H^{2+\epsilon}(\mathbb{R}^3)$ 中证明系统 (64) 的适定性. 为稍加简化, 考虑形式
\[
G^{\alpha\beta}(\phi)\partial_{\alpha}\partial_{\beta}\phi=N(\phi,\partial\phi)
\]
的波动方程, 满足初值条件 $\phi[0] \in H^s$. 这里 $G^{\alpha\beta}(\phi)$ 是 $\mathbb{R}^{1+3}$ 中一族 Lorentz 度量, 在原点邻域光滑依赖于 $\phi$, 而 $N$ 在 $\partial\phi$ 中为二次. 我们想证明对 $s > 2 + \epsilon$ 该 IVP 适定. 对形如
\[
\Box\phi=N(\phi,\partial\phi)
\]
的半线性波动方程的类似结果极易证明$^{21}$, 它只依赖于如下形式的 Strichartz 型估计:
\[
\|\partial\psi\|_{L_{t}^{2}L_{x}^{\infty}([0,T]\times\mathbb{R}^{n})}\lesssim\|\phi[0]\|_{H^{s}(\mathbb{R}^{n})}
\]
对任意 $s > \frac{n}{2} + \frac{1}{2}$、$n = 3$ 成立. 特别地它不要求非线性项满足如零条件的任何特殊结构.
\section*{3.9 拟线性波动方程的最优适定性}

我打算报告我与 I. Rodnianski \cite{klai42} 关于拟线性波动方程最优适定性问题的近期结果. 受形如 (64) 与 (67) 的拟线性波动方程启发, 我们考虑 $\mathbb{R}^{n+1}$($n \geq 3$)中的简化模型方程
\[
\partial_{t}^{2}\phi-g^{i j}(\phi)\partial_{i}\partial_{j}\phi=N(\phi,\partial\phi),
\]
其中 $N$ 为在 $\partial\phi$ 中二次的光滑函数. 基于 $H^{s}$ 能量估计与 Sobolev 嵌入定理的经典局部存在唯一性结果表明, (67) 在 Sobolev 空间 $H^{s}$(任意 $s > \frac{n}{2} + 1$)中局部适定. 由标准分部积分与交换子引理得到的 $H^{s}$ 能量估计取形式
\[
\|\phi(t)\|_{H^{s}}\leq M(\|\phi_{0}\|_{H^{s}}+\|\phi_{1}\|_{H^{s-1}}),
\]
其中 $M$ 连续依赖于积分 $\int_0^t \|\partial\phi\|_{L_x^\infty}$. 对充分小的时刻 $t$, Sobolev 估计 $\|\partial\phi\|_{L^\infty} \leq c\|\phi\|_{H^s}$($s > \frac{n}{2} + 1$)使我们能把上式右端纯用 $\|\phi_0\|_{H^s} + \|\phi_1\|_{H^{s-1}}$ 界定. 这些对 $s > \frac{n}{2} + 1$ 的局部先验估计可轻易化为局部存在唯一性结果.
量 $\int_0^t \|\partial\phi\|_{L_\infty^\infty}$ 可借助 Strichartz 估计 (66) 更好地控制. 这就带来了上述 $\frac{1}{2}$ 导数的获得.
在拟线性情形, 度量 $g^{ij}$ 依赖于解 $\phi$, 因而其正则性至多与 $\phi$ 相同. 因此需处理对线性方程 $\Box_g\phi = \partial_t^2 - g^{ij}\partial_t\partial_j\phi = 0$ 的 Strichartz 估计的证明, 其中系数与线性化波动算子的系数一致. 线性方程 $\Box_g$ 的系数所需条件须与 (67) 中非线性系数 $g^{ij}(\phi)$ 对 $\phi$ 的依赖一致. 鉴于 $\|\phi\|_{H^s}$ 的 Sobolev 范数应有界, 我们可设$^{22}$ $\|g\|_{H^s}$ 有界. 而且由于拟线性问题 (67) 在时刻 $t$ 的解的 $H^s$ 范数被初数据的 $H^s$ 范数控制, 只要 $\int_0^t\|\partial\phi\|_{L_\infty^\infty}$ 有界, 我们可归纳地假设度量 $g^{ij} = g^{ij}(\phi)$ 也满足 $\int_0^t\|\partial g\|_{L_\infty^\infty} \leq B_0$. 依我们在平坦空间 Strichartz 估计的经验, 我们不期望直接触及 $\|\partial\phi\|_{L_t^1L_\infty^\infty}$ 的范数, 但在局部时间它可由 $\|\partial\phi\|_{L_t^2L_\infty^\infty}$ 控制. 因此, 为闭合论证并导出 (67) 改进的局部存在唯一性结果, 我们需要证明 $H^s$ 范数的有界性, 以及对线性方程 $\Box_g\phi = 0$ 解的局部 $L_t^2L_\infty^\infty$-Strichartz 估计.
历史注记: 关于对系数粗糙$^{23}$ 的 $\Box_g \phi = 0$ 解之 Strichartz 型估计的首要重要工作归于 H. Smith \cite{klai60}. 他证明了标准 Strichartz 估计对 $C^{1,1}$ 系数方程成立. 但 $C^{1,1}$ 条件对应用于拟线性方程过强. 而且 H. Smith 与 C. Sogge \cite{klai61} 的一些重要反例表明, 若系数比 $C^{1,1}$ 更粗糙, 标准 Strichartz 估计失效.
这是 Bahouri 与 J.-Y. Chemin 的重要突破之前的状况. 在 \cite{klai1} 中他们证明, 对某些比 $C^{1,1}$ 更粗糙、与拟线性方程应用相容的系数, Strichartz 估计的较弱形式仍然存在. 即他们为变系数波动方程 $\Box_g \psi = 0$(系数满足上述类型假设)建立了带损失的 Strichartz 估计
\[
\left\|\partial\psi\right\|_{L_{t}^{2}L_{x}^{\infty}}\lesssim\left\|\psi[0]\right\|_{H^{\frac{n}{2}+\frac{1}{2}+\sigma}},
\]
只要损失 $\sigma \leq \frac{1}{2}$, 此估计可用于证明此处拟线性问题的局部适定性结果, 强于经典结果. H. Bahouri 与 J.-Y. Chemin 的结果基于对线性化问题 $\Box_g \psi = 0$(度量 $g$ 满足上述类型条件)带损失 $\sigma > \frac{1}{4}$ 的 Strichartz 估计. 后来 D. Tataru \cite{klai72} 证明 Strichartz 估计中的损失$^{24}$ 可改进为 $\sigma > \frac{1}{6}$. 事实上 \cite{klai72} 给出了低于 $C^{1,1}$ 时度量的光滑性与 Strichartz 估计中相应损失之间的精确关系. 最近 H. Smith 与 D. Tataru \cite{klai62} 证明这些结果也是尖锐的, 因此仅基于线性理论无法改进 Tataru 的结果.
Bahouri--Chemin 与 Tataru 的结果都基于两个主要要素.
a) 仿微分法: 考虑对应于空间频率 $\frac{\lambda}{2} \leq |\xi| \leq 2\lambda$($\lambda \in 2^\mathbf{N}$)截断的二进投影算子 $P_\lambda$. 令 $\phi^\lambda = \phi_\lambda = P_\lambda \phi$, 有 $\phi = \sum_\lambda \phi_\lambda$. 若 $\phi$ 满足拟线性波动方程 (67), 则 $\phi^\lambda$ 满足方程
\[
\Box_{g_{\leq\lambda^{a}}}\phi^{\lambda}=-\partial_{t}^{2}\phi^{\lambda}+g_{\leq\lambda^{a}}^{i j}\partial_{i}\partial_{j}\phi^{\lambda}=R_{\lambda}^{a},
\]
其中 $g_{\leq\lambda^a} = \sum_{\mu\leq\lambda^a} P_\mu g$, $0 < a < 1$, $R_\lambda$ 为易控制的误差项. Bahouri--Chemin 后来由 Tataru 精炼的原始思想, 是在频率依赖的时间区间 $I_{\lambda}$(大小 $|I_{\lambda}| \lesssim \lambda^{1-a}$)上, 对 (69) 证明无损失的 Strichartz 估计. 正是控制时间区间 $I_{\lambda}$ 大小上的这一损失, 导致了 (71) 中可微性的损失.
b) 基于参数化的色散不等式证明: 如我们所知, Strichartz 不等式证明的主要困难构件是能量与色散不等式. Bahouri--Chemin 与 Tataru 都利用参数化构造$^{25}$ 给出近似线性化方程 (69) 的解. 他们借助驻相法, 由参数化构造给出的近似解显式形式导出色散不等式.
c) 程函方程(eikonal equation): 两种方法的主要要素之一是构造程函方程的解
\[
u_{t}^{2}-g_{\leq\lambda^{a}}^{i j}\partial_{i}u\partial_{j}u=0,\qquad u(0,x,\xi)=x\cdot\xi,
\]
以及相伴的输运方程.
在 \cite{klai30} 中我发展了另一种更几何的方法, 基于第一讲讨论的向量场法的修改来证明 (69) 的色散不等式. 回忆命题 (1.9) 可知, 我们无需解的显式表示, 只用了标准波动方程与伸缩、推进及角动量向量场的交换性质, 就能证明色散不等式. 为对 (69) 型拟线性方程实施此策略, 我们需要类似的一族 Killing 与保形 Killing 向量场, 但对应度量 $g \leq \lambda^a$. 这当然不可能; 一般度量没有任何对称性. 我们能做的最好是构造形变张量尽可能小的向量场. 在 \cite{klai30} 中我概述了如何基于同一程函方程 (79) 但具不同初值(对应于此前的一族零超平面, 而此处对应于时间轴上顶点的一族出特征锥)来作此构造. 这一直接几何方法使我恢复了三维拟线性波动方程的 Tataru 结果.
利用 \cite{klai30} 中类似的几何方法, 最近 Rodnianski 与我在 \cite{klai42} 中, 通过考虑方程 (67) 的非线性特征, 改进了 Bahouri--Chemin 与 Tataru 的适定性结果. 我们不只对满足 $\|g\|_H^s$ 与 $\|\partial g\|_{L_1^1 L_\infty^s}$ 估计的 $\Box_g \phi = 0$ 解证明 Strichartz 估计, 而是本质地利用了方程 (67) 的系数 $g^{ij}(\phi)$ 自身满足形如 $\Box_g g_{ij} = N$(其中 $N$ 只依赖于 $\phi$ 与 $\partial \phi$)的方程这一事实.
我们的主要结果包含在如下定理中.
\setcounter{mtthm}{9}
\begin{theorem}\label{thm:3.10}
$\mathbb{R}^{3+1}$ 中的拟线性初值问题 (67) 在 $H^{s_*}$ 中局部适定, 其中 $s_* > s_0 = 2 + \frac{2 - \sqrt{3}}{2}$. 即对任意初值 $\phi[0] \in H^{s_*}$, 存在充分小的区间 $[0, T]$($T = T(\|\phi[0]\|_{H^{s_*}})$), 使得问题 (67) 有唯一解 $\phi \in C([0, T], H^{s_*})\cap C^{1}([0, T], H^{s_0-1})$. 此外 $\phi$ 满足 Strichartz 型估计
\[
\|\partial\phi\|_{L_{[0,T]}^{2}L_{x}^{\infty}}\leq c\, T^{s_{*}-s_{0}}\|\phi[0]\|_{H^{s_{*}}}.
\]
而且我们相信, 对 \cite{klai42} 的证明作相对较小修改, 我们现在已具备证明爱因斯坦方程$^{26}$ 在 $H^{2+\epsilon}$ 中局部适定所需的全部工具.\end{theorem}

\section*{3.11 \thmref{thm:3.10} 证明概要}

经过一系列如今已典范的简化(见本讲补充), 我们可把定理的证明化为如下情形:
设给定一族$^{27}$ 黎曼度量 $h_{ij}^{\lambda}(t,x)dx^{i}dx^{j}$, 定义在时间板 $I_{\lambda} \times \mathbb{R}^{3}$ 上, $I_{\lambda} = [0, t_{*}]$, $t_{*} \leq \lambda^{a}$, 满足如下条件
\[
\|\partial^{1+m}h_{\lambda}\|_{L_{I_{\lambda}}^{1}L_{x}^{\infty}}\lesssim\lambda^{-(1-a)(m+1)},\qquad
\|\partial^{1+m}h_{\lambda}\|_{L_{I_{\lambda}}^{2}L_{x}^{\infty}}\lesssim\lambda^{-\frac{2-a}{2}-(1-a)m},
\]
\[
\|\partial^{1+m}h_{\lambda}\|_{L_{I_{\lambda}}^{\infty}L_{x}^{\infty}}\lesssim\lambda^{-1+\frac{a^{2}}{2}-(1-a)m},\qquad
\|\partial_{x}^{\frac{1}{2}+m}(\partial^{2}h_{\lambda})\|_{L_{I_{\lambda}}^{\infty}L_{x}^{2}}\lesssim\lambda^{-1+\frac{a^{2}}{2}-(1-a)m},
\]
\[
\|\partial^{m}\bar{\Box}_{h_{\lambda}}h_{\lambda}\|_{L_{I_{\lambda}}^{1}L_{x}^{\infty}}\lesssim\lambda^{-(2-a)-(1-a)m},
\]
其中 $\bar{\Box}_{h_{\lambda}} = \partial_{t}^{2} - h_{\lambda}^{ij} \partial_{i} \partial_{j}$. 我们需证明对协变波动方程
\[
\Box_{h_{\lambda}}\psi=0
\]
的解的如下衰减估计, 其中
\[
\Box_{h_{\lambda}}\psi=\frac{1}{\sqrt{\det h_{\lambda}}}\partial_{t}\sqrt{\det h_{\lambda}}\partial_{t}\psi-\frac{1}{\sqrt{\det h_{\lambda}}}\partial_{i}\bigl(h_{\lambda}^{i j}\sqrt{\det h_{\lambda}}\partial_{j}\psi\bigr).
\]

\setcounter{mtthm}{11}
\begin{theorem}\label{thm:3.12}
设 $\psi$ 为 $[0, t_*] \times \mathbb{R}^n$ 中 (77) 的解, $t_* \leq \lambda^a$, 其初数据 $\psi[0]$ 支在物理空间中以原点为中心、半径 $\frac{1}{2}$ 的球内. 设度量 $h_\lambda$ 满足 (138)--(142) 且 $a < -1 + \sqrt{3}$. 则对所有 $t \in [0, t_*]$,
\[
\|\partial\psi(t)\|_{L^{\infty}}\leq(1+t)^{-1+\epsilon}\|\psi[0]\|_{H^{k}}.
\]
\end{theorem}

证明基于第三讲补充中讨论的\propref{prop:4.11} 的弯曲时空类似. 在陈述它之前我们需要一些定义:
\setcounter{mtthm}{12}
\begin{definition}\label{def:3.13}
零对(null pair)由两向量场 $L,\underline{L}$ 组成, 满足
\[
H(L,L)=H(\underline{L},\underline{L})=0,\qquad H(L,\underline{L})=-2.
\]
与零对 $L, \underline{L}$ 相伴的零标架(null frame)由 4 个线性无关向量 $e_1, e_2, e_3 = \underline{L}, e_4 = L$ 组成, 满足
\[
H(e_{A},L)=H(e_{A},\underline{L})=0,\quad H(e_{A},e_{B})=\delta_{A B},\quad\text{对 }A=1,2.
\]
\end{definition}

\setcounter{mtthm}{13}
\begin{definition}\label{def:3.14}
光学函数$^{28}$ $u$ 定义为程函方程
\[
H^{\alpha\beta}\partial_{\alpha}u\partial_{\beta}u=0,\qquad\text{或}\quad\partial_{t}u=|\nabla u|_{h}
\]
的向前解, 边界条件为在时间轴上 $u = t$. 时间函数 $t$ 的等值超曲面 $\Sigma_t$ 与 $u$ 的等值超曲面 $C_u$ 的交的 2-曲面记为 $S_{t,u}$.\end{definition}

\setcounter{mtthm}{14}
\begin{definition}\label{def:3.15}
给定这样的 $u$, 我们定义典范零对如下:
\[
L=T+N,\qquad \underline{L}=T-N,
\]
其中
\[
T=\partial_{t},\qquad N=-\frac{h^{i j}\partial_{j}u}{|\nabla u|_{h}}\partial_{i}.
\]
注意 $T + N$ 是零锥 $C_u$ 的生成法向量. 我们可补充 $L, \underline{L}$ 构成张成整个切空间的零标架. 这可通过在 $S_{t,u} = C_u \cap \Sigma_t$ 上任取一组正交标架 $e_A$ 实现. 我们重写完整零标架 $e_1, e_2, e_3 = \underline{L}, e_4 = L$.
我们还引入如下:
\begin{enumerate}[label=\roman*)]
\item 能量--动量张量: $Q_{\alpha\beta} = \partial_{\alpha}\psi\partial_{\beta}\psi - \tfrac{1}{2}H_{\alpha\beta}(H^{\mu\nu}\partial_{\mu}\psi\partial_{\nu}\psi)$.
\item 修正能量--动量张量: $\bar{Q}_{\alpha\beta} = Q_{\alpha\beta} + (n-1)t\psi\partial_{\beta}\psi - \frac{n-1}{2}\psi^{2}$.
\item 修正 Morawetz 向量场: 令 $K = \tfrac{1}{2}(\underline{u}^2L + u^2\underline{L})$, 其中 $\underline{u} = 2t - u$.
\item 修正形变张量: $\bar{\pi} = (K)\pi - 4tH$.
\item 共形能量: $Q_0[\psi](t) = \int_{\Sigma_t} \bar{Q}(K, \partial_t)[\psi]$.
\item 全共形能量: $Q[\psi](t) = Q_0[\psi](t) + Q_0[\partial_t\psi](t) + Q_0[\partial_t^2\psi](t)$.
\item $Q$ 的零分量: 相对零标架 $Q$ 的分量为
\[
Q_{44}=(D_{4}\psi)^{2},\quad Q_{34}=|\bar{\nabla}\psi|^{2},\quad
Q_{33}=(D_{3}\psi)^{2},\quad Q_{3A}=D_{3}\psi\bar{\nabla}_{A}\psi,
\]
\[
Q_{4A}=D_{4}\psi\bar{\nabla}_{A}\psi,\qquad Q_{A B}=\bar{\nabla}_{A}\psi\bar{\nabla}_{B}\psi+\frac{1}{2}(D_{3}\psi D_{4}\psi-|\bar{\nabla}\psi|^{2})\delta_{A B}.
\]
\item 共形能量密度: 由 vii) 易计算:
\[
Q(K,\partial_{t})[\psi]=\frac{1}{4}\bigl(\underline{u}^{2}(D_{4}\psi)^{2}+u^{2}(D_{3}\psi)^{2}+(\underline{u}^{2}+u^{2})|\nabla\psi|^{2}\bigr).
\]
\end{enumerate}
\end{definition}

\thmref{thm:3.12} 的证明可化为如下:
\setcounter{mtthm}{15}
\begin{theorem}[有界性定理]\label{thm:3.16}
在\thmref{thm:3.12} 同样假设下, 有
\[
\mathcal{Q}(t)\lesssim\mathcal{Q}(0).
\]
\end{theorem}

下面我们勾画有界性定理证明的一些主要思想. 我们限于讨论 $Q_{0}$ 的有界性.
步骤 1. 引入量
\[
\mathcal{E}_{0}[\psi](t)=\int_{\Sigma_{t}}\bigl(\underline{u}^{2}(D_{4}\psi)^{2}+u^{2}(D_{3}\psi)^{2}+\underline{u}^{2}|\bar{\nabla}\psi|^{2}+\psi^{2}\bigr)\lesssim\mathcal{Q}_{0}[\psi](t).
\]
由类似于\propref{prop:4.11} 之 iii) 的比较定理, 我们证明
\[
\mathcal{E}_{0}[\psi](t)\lesssim\mathcal{Q}_{0}[\psi](t).
\]

步骤 2. \thmref{thm:3.16} 证明的关键要素是如下\propref{prop:4.11} 的广义共形恒等式的弯曲时空版本.
\setcounter{mtthm}{16}
\begin{proposition}[广义共形恒等式]\label{prop:3.17}
设 $\psi$ 为 $\Box_h \psi = 0$ 的解. 则
\[
\mathcal{Q}_{0}[\psi](t)=\mathcal{Q}_{0}[\psi](t_{0})-\frac{1}{2}\int_{[t_{0},t]\times\mathbb{R}^{3}}Q^{\alpha\beta}\bar{\pi}_{\alpha\beta}\,d\omega+\frac{n-1}{2}\int_{[t_{0},t]\times\mathbb{R}^{3}}\psi^{2}\Box_{h}(t)\,d\omega.
\]
\end{proposition}

步骤 3. 为证明有界性定理, 需控制 $\int_{[t_0,t]\times\mathbb{R}^3}Q^{\alpha\beta}\bar{\pi}_{\alpha\beta}$. 相对零标架展开, 有
\[
Q^{\alpha\beta}\bar{\pi}_{\alpha\beta}
=\frac{1}{4}Q_{44}\bar{\pi}_{33}+\frac{1}{4}Q_{33}\bar{\pi}_{44}+\frac{1}{2}Q_{34}\bar{\pi}_{34}
-Q_{3A}\bar{\pi}_{4A}-Q_{4A}\bar{\pi}_{3A}+Q_{A B}\bar{\pi}_{A B}.
\]
我们需要形如
\begin{align*}
\int_{[t_{0},t]\times\mathbb{R}^{3}}Q^{\alpha\beta}\bar{\pi}_{\alpha\beta}
&\lesssim\lambda^{-\epsilon}\sup_{[t_{0},t]}\int_{\Sigma_{\tau}}\bigl(\underline{u}^{2}(D_{4}\psi)^{2}+u^{2}(D_{3}\psi)^{2}+\underline{u}^{2}|\nabla\psi|^{2}+\psi^{2}\bigr)\\
&=\lambda^{-\epsilon}\sup_{[t_{0},t]}\mathcal{E}_{0}[\psi](\tau)\lesssim\lambda^{-\epsilon}\sup_{[t_{0},t]}\mathcal{Q}_{0}[\psi](\tau)
\end{align*}
的估计, 对任意 $\epsilon > 0$. 特别考虑项
\[
\int_{[t_{0},t]\times\mathbb{R}^{3}}Q_{33}\bar{\pi}_{44}=\int_{[t_{0},t]\times\mathbb{R}^{3}}\left|D_{3}\psi\right|^{2}\bar{\pi}_{44}.
\]
为得到所需估计, 我们需要
\[
|\bar{\pi}_{44}(\tau,\cdot)|\lesssim u^{2}\tau^{-1}\lambda^{-\epsilon}.
\]
换言之在波区 $|u| << t$ 外 $\bar{\pi}_{44}$ 必须很小. 事实上为估计 $\int Q^{\alpha\beta}\bar{\pi}_{\alpha\beta}$ 中多数项, 我们需要对 $K$ 的修正形变张量 $\bar{\pi}$ 的如下估计.
\setcounter{mtthm}{17}
\begin{proposition}\label{prop:3.18}
$K$ 的修正形变张量的零分量满足估计:
\[
|\bar{\pi}_{44}|\leq u^{2}\lambda^{-\epsilon}\tau^{-1},\qquad|\bar{\pi}_{34}|\leq\underline{u}^{2}\lambda^{-\epsilon}\tau^{-1},
\]
\[
|\bar{\pi}_{33}|\leq\underline{u}^{2}\lambda^{-\epsilon}\tau^{-1},\quad|\bar{\pi}_{3A}|\leq\underline{u}^{2}\lambda^{-\epsilon}\tau^{-1},
\]
\[
|\bar{\pi}_{4A}|\leq u\underline{u}\lambda^{-\epsilon}\tau^{-1},\quad|\bar{\pi}_{A B}|\leq\underline{u}^{2}\lambda^{-\epsilon}\tau^{-1}.
\]
\end{proposition}

步骤 4. 上述估计可处理除项
\[
\frac{1}{2}\int_{[t_{0},t]\times\mathbb{R}^{3}}\bar{\pi}D_{3}\psi D_{4}\psi,\qquad\text{其中}\quad \bar{\pi}=\delta^{A B}\bar{\pi}_{A B}
\]
外的所有项. 为估计此项我们需要分部积分论证, 它将引入 $D_4\bar{\pi}$ 与 $\bar{\nabla}\bar{\pi}$, 对它们我们需要类似于\propref{prop:3.18} 的估计.
步骤 5. 我们需要证明\propref{prop:3.18} 的渐近估计, 以及对 $D_4\bar{\pi}$ 与 $\bar{\nabla}\pi$ 的类似估计. 这依赖于对程函方程 (79) 的关键分析. 事实上易见, 由直接计算 $\bar{\pi}$ 依赖于时间叶层 $\Sigma_t$ 的第 II 基本形式 $k$ 以及更重要的 $u$ 的 Hessian $D^2u$. 由于 $k_{ij} = -\frac{1}{2}\partial_t h_{ij}$, 我们对 $k$ 有很好控制. 另一方面 $u$ 的 Hessian 满足 Ricatti 型方程
\[
D_{L}(D^{2}u)+D^{2}u\cdot D^{2}u=R(L,L),
\]
其中 $R(L,L)$ 为曲率张量 $R_{\alpha\beta\gamma\delta}$ 与 $L$ 二次缩并的分量. 初看似乎 $D^{2}u$ 具有与曲率张量 $R$ 相同的可微性. 而且为在某时刻 $t$ 估计 $D^{2}u$, 需在零锥 $C_{u}$ 的测地线生成元上积分 $R(L,L)$. 因此, 由假设 (138), 由于 $R$ 依赖于 $h$ 的两个导数, 我们似乎能得到的最佳估计是
\[
\|D^{2}u(t)\|_{L^{\infty}}
\lesssim\int_{0}^{t}|\partial^{2}h(\tau,\cdot)|\,d\tau
\lesssim\lambda^{-2(1-a)}.
\]
沿此路线的论证只会给出 Tataru 得到的最优结果. 要做得更好我们必须利用非线性假设 (142). 能利用该假设是我们几何方法的关键优势.
为此我们需要相对典范零标架分解 $D^{2}u$. 这可通过引入如下量实现:
\[
\angle_{A B}=\langle D_{A}e_{4},e_{B}\rangle,\qquad \eta_{A}=\frac{1}{2}\langle D_{3}e_{4},e_{A}\rangle.
\]
又令
\[
b^{-1}=\partial_{t}u=|\nabla u|_{h}.
\]
我们把 $\chi_{AB}$ 分解为 $\mathrm{tr}\chi = \delta^{AB}\chi_{AB}$ 与无迹部分 $\hat{\chi}_{AB}$. 它们满足方程:
\[
L(b)=-b\, k_{N N},
\]
\[
L(\mathrm{tr}\chi)+\frac{1}{2}(\mathrm{tr}\chi)^{2}=-|\hat{\chi}|^{2}-k_{N N}\mathrm{tr}\chi-R_{44},
\]
\[
\mathcal{P}_{4}\hat{\chi}_{A B}+\frac{1}{2}(\mathrm{tr}\chi)\hat{\chi}_{A B}=-k_{N N}\hat{\chi}_{A B}-\hat{\alpha}_{A B},
\]
\[
\mathcal{P}_{4}\eta_{A}+\frac{1}{2}(\mathrm{tr}\chi)\eta_{A}=-(k_{B N}+\eta_{B})\hat{\chi}_{A B}-\frac{1}{2}(\mathrm{tr}\chi)k_{A N}-\frac{1}{2}\beta_{A},
\]
其中 $\alpha_{AB} = R_{A4B4}$, $\hat{\alpha} = \delta^{AB} \alpha_{AB}$, $\beta_{A} = R_{A43B}$. 若我们试图从各自输运方程 (89)、(90) 与 (91) 估计 $\mathrm{tr}\chi$、$\hat{\chi}$、$\eta$, 会遇到上面提到的同样困难. 我们用如下两个思想绕过这些困难:
思想 1. 注意出现在 $\mathrm{tr}\chi$ 输运方程中的曲率项 $R_{44}$ 有特殊结构. 事实上我们有

\setcounter{mtthm}{18}
\begin{lemma}[显著分解]\label{lem:3.19}
$R_{44} = e_4^\mu e_4^\nu R_{\mu\nu}$(Ricci 曲率的)分量容许如下分解:
\[
R_{44}=L(z)-\frac{1}{2}e_{4}^{\mu}e_{4}^{\nu}\Box_{h}H_{\mu\nu}+Error,
\]
其中函数 $z$ 与 Error 满足点态估计 $|z| \lesssim |\partial H|$, $|Error| \lesssim (\partial H)^2$.\end{lemma}

由该引理我们可把 $\operatorname{tr}\chi$ 的输运方程改写为:
\setcounter{mtthm}{19}
\begin{proposition}\label{prop:3.20}
令 $y = \operatorname{tr}\chi - \tfrac{2}{s}$, 其中 $s = t - u$, $z$ 与 error 如上面引理. 则
\[
L(y+z)+\operatorname{tr}\chi(y+z)
=\frac{1}{2}(y+z)^{2}+\frac{2}{s}z-\frac{1}{2}z^{2}-|\hat{\chi}|^{2}
-k_{N N}\operatorname{tr}\chi+\frac{1}{2}e_{4}^{\nu}e_{4}^{\mu}\Box_{h}H_{\mu\nu}-Error.
\]
\end{proposition}

思想 2: $y + z$ 的方程看来比 $\text{tr}\chi$ 的方程好得多. 不幸的是对 $\hat{\chi}$ 与 $\eta$ 不能做同样技巧. 而且 $y + z$ 的方程本身与 $\hat{\chi}$ 的输运方程耦合. 为解决这个问题我们求助于一个完全不同的几何方程, 即 $\chi$ 的 Codazzi 方程. 它有形式
\[
(\mathrm{div}\,\hat{\chi})_{A}+\hat{\chi}_{A B}k_{B N}=\frac{1}{2}(\nabla_{A}\mathrm{tr}\chi+k_{A N}\mathrm{tr}\chi)-R_{B4A B}.
\]
注意方程 $\mathrm{div}\,\hat{\chi}_{A}=F_{A}$ 是曲面 $S_{t,u}$ 上的椭圆 Hodge 系统. 利用上述 $y+z$ 的输运方程与 (93) 的耦合系统, 我们导出形如
\[
|\mathrm{tr}\chi-\frac{2}{s}|\lesssim\lambda^{-\bar{a}},\qquad |\hat{\chi}|\lesssim\lambda^{-\bar{a}}
\]
的估计, 其中
\[
\bar{a}=1-\frac{a^{2}}{2}.
\]
这远好于前面得到的 $\lambda^{-2(1-a)}$ 估计.
\setcounter{mtthm}{20}
\begin{remark}\label{rem:3.21}
对 $\eta$ 我们有类似方法. 思想是得到 $\eta$ 的散度--旋度系统, 耦合于一个类似于 $\text{tr}\chi$ 的输运方程.\end{remark}

\setcounter{mtthm}{21}
\begin{remark}\label{rem:3.22}
利用耦合椭圆 Hodge 系统与有利输运方程的思想, 出现在非常不同背景的 \cite{klai12} 中, 与 Minkowski 空间的全局稳定性有关.\end{remark}

\chapter{第一讲补编}

重要的非线性波动方程, 如 Yang--Mills 方程、波映射与 Einstein 方程, 其基本对称性都与其几何特征紧密相联. 在研习这些方程时, 微分几何(特别是 Lorentz 几何)的知识, 与调和分析至少同等重要. 下面我们将考察 Minkowski 时空的共形结构与该结构同波动方程 (1) 之估计之间的关系. 特别地, 我们将用到第三讲中讨论的广义能量估计. 本文假定读者熟悉微分几何的基本概念, 如 Lie 微分、协变微分以及曲率张量.
\section*{4.1 Minkowski 时空的共形结构}

\setcounter{mtthm}{1}% 4.1 为节标题, 首个定理环境显示 4.2

\begin{definition}\label{def:4.2}
考虑一个时空 $(\mathcal{M}, g)$, 即一个 Lorentz 流形, 其中 $g$ 为号差 $(-1,1,\ldots,1)$ 的 Lorentz 度量.
一个微分同胚 $\Phi:\mathcal{U}\subset\mathcal{M}\to\mathcal{M}$ 称为\textbf{共形等距}, 如果对任意点 $p$ 均有 $\Phi_* g=\Lambda^2 g$, 即
\begin{equation}
(\Phi^* g)(X,Y)|_p=g(\Phi_* X,\Phi_* Y)|_{\Phi(p)}=\Lambda^2 g(X,Y)|_p
\end{equation}
且 $\Lambda\neq 0$. 若 $\Lambda=1$, 则称 $\Phi$ 为 $\mathcal{M}$ 的\textbf{等距}.\end{definition}

\begin{definition}\label{def:4.3}
若一个向量场 $K$ 生成一参数等距群(相应地, 共形等距群), 则称 $K$ 为\textbf{Killing 向量场}(相应地, \textbf{共形 Killing 向量场}).\end{definition}

设 $K$ 为这样的向量场, $\Phi_t$ 为相应的单参数群. 由于 $(\Phi_t)_*$ 均为共形等距, 可知 $\mathcal{L}_K g$ 必与度量 $g$ 成比例; 若 $K$ 为 Killing 向量场, 则 $\mathcal{L}_K g=0$.
注意, C(共形结构)关于共形等距具有共形不变性这一重要性质.
\begin{definition}\label{def:4.4}
对任意向量场 $X$, 记 $(X)_\pi$ 为 $X$ 的\textbf{形变张量}, 由下式定义:
\begin{equation}
{}^{(X)}\pi_{\alpha\beta}=(\mathcal{L}_X g)_{\alpha\beta}=D_\alpha X_\beta+D_\beta X_\alpha.\end{equation}
\end{definition}

该张量 $(X)\pi$ 以一种精确的意义, 衡量了由 $X$ 生成的微分同胚偏离等距或共形等距的程度.
\begin{proposition}\label{prop:4.5}
向量场 $X$ 为 Killing 向量场, 当且仅当 $(X)\pi=0$; 为共形 Killing 向量场, 当且仅当 $(X)\pi$ 与 $g$ 成比例.\end{proposition}

为证明"当"的部分, 可取局部坐标 $x^0,x^1,\ldots,x^n$ 使得 $X=\frac{\partial}{\partial x^0}$. 由此立即可见, 相对这些坐标, 度量 $g$ 与 $x^0$ 无关.
\begin{proposition}\label{prop:4.6}
任意时空 $\mathcal{M}$ 上, 线性无关的 Killing 向量场个数至多为 $\frac12(n+1)(n+2)$.\end{proposition}

拥有最多 Killing 与共形 Killing 向量场的时空正是 Minkowski 时空. 下面我们复习其相应的等距与共形等距.
设 $x^\mu$ 为正定向的惯性坐标系, 我们有:
\begin{enumerate}
\item 平移: 对任意给定向量 $a=(a^0,a^1,\ldots,a^n)\in\mathbf{M}$,
\begin{equation}
x^\mu\to x^\mu+a^\mu.\end{equation}
\item Lorentz 旋转: 对任意 $\Lambda=\Lambda^\rho_\sigma\in\mathbf{O}(1,n)$,
\begin{equation}
x^\mu\to \Lambda^\mu_\nu x^\nu.\end{equation}
\item 伸缩: 对任意实数 $\lambda\neq 0$,
\begin{equation}
x^\mu\to \lambda x^\mu.\end{equation}
\item 反演: 考虑变换 $x^\mu\to I(x^\mu)$, 其中
\begin{equation}
I(x^\mu)=\frac{x^\mu}{(x,x)},\end{equation}
对一切满足 $(x,x)\neq 0$ 的点 $x\in\mathbf{M}$ 定义.\end{enumerate}

前两类变换是 $\mathbf{M}$ 的等距, 由它们生成的群称为 Poincar\'e 群; 后两类变换是共形等距, 由上述全部变换生成的群称为共形群. 下面列出生成上述变换的 Killing 与共形 Killing 向量场.
\begin{enumerate}[label=\roman*.]
\item $x^\mu$ 方向的平移生成元, $\mu=0,1,\ldots,n$:
\begin{equation}
T_\mu=\frac{\partial}{\partial x^\mu}.\end{equation}
\item $(\mu,\nu)$ 平面中 Lorentz 旋转的生成元:
\begin{equation}
L_{\mu\nu}=x_\mu\partial_\nu-x_\nu\partial_\mu.\end{equation}
\item 伸缩变换的生成元:
\begin{equation}
S=x^\mu\partial_\mu.\end{equation}
\item 反演平移的生成元:
\begin{equation}
K_\mu=2x_\mu\Bigl(x^\rho\frac{\partial}{\partial x^\rho}-(x^\rho x_\rho)\frac{\partial}{\partial x^\mu}\Bigr).\end{equation}
\end{enumerate}

\begin{theorem}\label{thm:4.7}
$\mathbf{M}^{n+1}$ 中任意共形 Killing 向量场($n>1$)都是向量场 $T,L,S,K$ 以实常系数为系数的线性组合. 若 $n=1$, 则所有共形向量场由下式给出:
\begin{equation}
f(t+x)(\partial_t+\partial_x)+g(t-x)(\partial_t-\partial_x),\end{equation}
其中 $f,g$ 为任意函数.\end{theorem}

\setcounter{mtthm}{8}% 4.8 为节标题, 下一定理环境显示 4.9

\section*{4.8 广义能量估计}

考虑平直 Minkowski 空间 $\mathbb{R}^{n+1}$ 中的波动方程 $\Box\phi=0$ 及其相伴的能量---动量张量
\begin{equation}
Q_{\alpha\beta}=Q_{\alpha\beta}[\phi]=\partial_\alpha\phi\partial_\beta\phi-\frac12 m_{\alpha\beta}\bigl(m^{\mu\nu}\partial_\mu\phi\partial_\nu\phi\bigr).\end{equation}

\begin{proposition}\label{prop:4.9}
能量---动量张量是对称且无散度的, 即
\begin{equation}
\partial^\beta Q_{\alpha\beta}=0.\end{equation}
此外, 对任意指向未来的类时向量场 $X,Y$, 均有 $Q(X,Y)\geq 0$.\end{proposition}

\begin{proof}
前两个结论显然. 为验证第三个, 在任一点取一对零向量 $L,\underline{L}$ 张成包含 $X,Y$ 的平面, 于是 $X,Y$ 都是 $L,\underline{L}$ 带正系数的线性组合. 由如下简单而重要的计算可知其正性:
\begin{gather}
Q(L,L)=L(\phi)^2,\\
Q(L,\underline{L})=|\bar\nabla\phi|^2,\\
Q(\underline{L},\underline{L})=L(\phi)^2,\end{gather}
其中 $|\bar\nabla\phi|=\sum_A |e_A(\phi)|^2$, $(e_A)_{A=1,\ldots,n-1}$ 为张成 $L,\underline{L}$ 正交补的标准正交标架.\end{proof}

设 $X$ 为任意向量场, 考虑与之相伴的时空动量向量场 $P_\alpha=Q_{\alpha\beta}X^\beta$. 利用 $Q$ 的对称性与无散度性质, 连同形变张量 $\pi=(X)\pi$ 的定义, 推得
\begin{equation}
\partial_\alpha P^\alpha=\frac12 Q^{\alpha\beta}\pi_{\alpha\beta}.\end{equation}

因此, 若 $X$ 为 Killing 向量场, 则 $\partial_\alpha P^\alpha=0$, 积分后即得熟知的守恒律. 特别地, 若我们在时空板 $[t_0,t]\times\mathbb{R}^n$ 上积分, 便得到
\begin{equation}
\int_{\Sigma_t}Q(X,\partial_t)\,dx=\int_{\Sigma_{t_0}}Q(X,\partial_t)\,dx,\end{equation}
其中 $\partial_t$ 在此扮演类空超曲面 $\Sigma_t$ 的未来指向单位法向. 取 $X$ 亦为 Killing 向量场 $\partial_t$, 即得熟知的能量守恒律
\begin{equation}
\int_{\Sigma_t}Q(\partial_t,\partial_t)\,dx=\int_{\Sigma_{t_0}}Q(\partial_t,\partial_t)\,dx,\end{equation}
其中 $Q(\partial_t,\partial_t)=\frac12\bigl(|\partial_t\phi|^2+\sum_{i=1}^n|\partial_i\phi|^2\bigr)$ 为通常的能量密度. 线性动量与角动量的守恒律也包含在上式中, 只需取 $X=\partial_i$ 或 $X=L_{ij}$.
若我们不在板 $[0,T]\times\mathbb{R}^n$ 上积分恒等式 $\partial^\alpha P_\alpha$, 而是在未来区域
\begin{equation}
\mathcal{D}=\{(t,x)\mid |x-y|\leq T-t,\; t_0\leq t\leq t_1\leq T\}
\end{equation}
上积分, 则得到守恒律
\begin{equation}
\int_{\Sigma_{t_1}}Q(X,\partial_t)\,dx+\int_{\underline{C}(t_0,t_1)}Q(X,\underline{L})=\int_{\Sigma_{t_0}}Q(X,\partial_t)\,dx,\end{equation}
其中 $\underline{C}(t_0,t_1)$ 为 $\mathcal{D}$ 的外边界, $\underline{L}=\partial_t-\sum_i\frac{(x^i-x_0^i)}{|x-x_0|}\partial_i$ 为指向入射零超曲面 $\underline{C}$ 的法向零向量场, 而通量 $\int_{\underline{C}(t_0,t_1)}Q(X,\underline{L})=\int_{t_0}^{t_1}dt\int_{S_{\underline{u}_t}}Q(X,\underline{L})$, 其中 $S_{\underline{u}_t}$ 为 $\underline{C}$ 与 $\Sigma_t$ 的交球面. 在 $X=\partial_t$ 的特殊情形, $ \Sigma_t $ 上的能量积分与通量均为正, 从而给出波动方程依赖域性质的简单证明.
上面导出 (97) 的程序可推广到共形 Killing 向量场.
\begin{proposition}\label{prop:4.10}
设 $Q_{\alpha\beta}=Q_{\alpha\beta}[\phi]$ 为 $\Box\phi=0$ 之解相伴的能量---动量张量. 令 $X$ 为共形 Killing 向量场, 即 $\pi=(X)\pi=\mathcal{L}_X m=\Omega m$, 且 $\mathrm{tr}\,\pi=m^{\alpha\beta}\pi_{\alpha\beta}$. 易验证 $\Box\Omega=0$; 特别地, 当 $X=K_0$ 时 $\Omega=4(n+1)t$. 记
\begin{equation}
\bar P_\alpha=Q_{\alpha\beta}X^\beta+\frac{n-1}{4(n+1)}\mathrm{tr}\,\pi\,\phi\partial_\alpha\phi-\frac{n-1}{8(n+1)}\partial_\alpha(\mathrm{tr}\,\pi)\phi^2.\end{equation}
则有
\begin{equation}
\partial^\alpha\bar P_\alpha=0.\end{equation}
\end{proposition}

将命题应用于共形 Killing 向量场 $X=K:=K_0$, 并在时间板 $[t_0,t]\times\mathbb{R}^n$ 上积分相应的无散度方程, 推得:
\begin{proposition}\label{prop:4.11}
设 $\overline{Q}(K_0,T_0)=Q(K_0,T_0)+(n-1)t\phi\partial_t\phi-\frac{n-1}{2}\phi^2$, 其中 $T_0=\partial_t$ 为 $\Sigma_t$ 的单位法向, $\phi$ 为 $\Box\phi=0$ 的解.
\begin{enumerate}[label=\roman*.]
\item 如下共形守恒律成立:
\begin{equation}
\int_{\Sigma_t}\overline{Q}(K_0,T_0)=\int_{\Sigma_{t_0}}\overline{Q}(K_0,T_0).\end{equation}
\item 此外,
\begin{equation}
\int_{\Sigma_t}\overline{Q}(K_0,T_0)=\frac14\Bigl(\int_{\Sigma_t}\underline{u}^2(L'\phi)^2+\int_{\Sigma_t}2(t^2+r^2)|\nabla\phi|^2+\int_{\Sigma_t}u^2(\underline{L}'\phi)^2\Bigr),\end{equation}
其中 $L=\partial_t+\partial_r$, $\underline{L}=\partial_t-\partial_r$, $u=t-r$, $\underline{u}=t+r$, 且 $\underline{u}L'(\phi)=\underline{u}L(\phi)+(n-1)\phi$, $u\underline{L}'(\phi)=u\underline{L}(\phi)+(n-1)\phi$.
\item 又若 $n\geq 3$, 则存在常数 $c>0$ 使得
\begin{equation}
\int_{\Sigma_t}\overline{Q}(K_0,T_0)\geq c\Bigl(\int_{\Sigma_t}\underline{u}^2(L\phi)^2+\int_{\Sigma_t}2(t^2+r^2)|\nabla\phi|^2+\int_{\Sigma_t}u^2(\underline{L}\phi)^2+\phi^2\Bigr).\end{equation}
\end{enumerate}
\end{proposition}

\begin{definition}\label{def:4.12}
我们引入如下范数:
\begin{gather}
E[\phi](t)=\int_{\Sigma_t}|\phi(t)|^2,\\
\mathcal{Q}[\phi](t)=\int_{\Sigma_t}\overline{Q}[\phi](K_0,T_0),\\
\mathcal{E}[\phi](t)=\int_{\Sigma_t}\Bigl(\underline{u}^2|L\phi|^2+2(t^2+r^2)|\nabla\phi|^2+u^2|\underline{L}\phi|^2+|\phi|^2\Bigr).\end{gather}
我们还引入高阶范数:
\begin{gather}
E_s[\phi](t)=\sum_{j=0}^s\;\sum_{\Gamma_{i_1},\ldots,\Gamma_{i_j}}E[\Gamma_{i_1},\ldots,\Gamma_{i_j}\phi](t),\\
\mathcal{Q}_s[\phi](t)=\sum_{j=0}^s\;\sum_{\Gamma_{i_1},\ldots,\Gamma_{i_j}}\mathcal{Q}[\Gamma_{i_1},\ldots,\Gamma_{i_j}\phi](t),\end{gather}
其中 $\Gamma$ 为任一向量场
\begin{equation}
S=t\partial_t+\sum_i x_i\partial_i,\qquad L_i=x_i\partial_t+t\partial_i,\qquad \Omega_{ij}=x_i\partial_j-x_j\partial_i.\end{equation}
\end{definition}

\setcounter{mtthm}{12}% 4.13 与 4.14 之间无跳过, 但 4.13 为注, 需钉号

\begin{remark}\label{rem:4.13}
注意
\begin{equation}
E_{s+1}(t)\lesssim\mathcal{Q}_s(t)=\mathcal{Q}_s(0).\end{equation}
\end{remark}

\setcounter{mtthm}{14}% 4.14 为节标题, 下一定理环境显示 4.15

\section*{4.14 衰减估计}

众所周知, 自由波动方程 $\Box\phi=0$ 在 $\mathbb{R}^{n+1}$ 中的解, 在一致范数下随 $t\to\infty$ 按 $t^{-\frac{n-1}{2}}$ 衰减. 此外, 远离出射特征方向时衰减更快. 这类信息可由物理空间或傅里叶空间中解的显式表示导出. 本节的目标是说明: 无需使用显式表示, 仅借助前一节讨论的广义能量估计, 便能获得此类信息.
考虑标准的零向量对 $E_\pm=\partial_t\pm\partial_r$, 以及角向量场 $A_i=\partial_i-\frac{x_i}{r}\partial_r$. 显然
\begin{equation}
\sum_i|A_i\phi|\lesssim|\bar\nabla\phi|\lesssim\sum_i|A_i\phi|.\end{equation}
又有
\begin{equation}
|\partial_r\phi|+\sum_i|A_i\phi|\lesssim|D\phi|\lesssim|\partial_r\phi|+\sum_i|A_i\phi|.\end{equation}
我们还可轻易验证如下简单的代数恒等式:
\begin{gather}
\frac12\bigl((t+r)E_+ +(t-r)E_-\bigr)=S,\\
\frac12\bigl((t+r)E_+ -(t-r)E_-\bigr)=\sum_i\frac{x_i}{|x|}L_i,\\
tA_i=L_i-\frac{x_i}{|x|}\sum_j\frac{x_j}{|x|}L_j,\\
t\Omega_{ij}=x_iL_j-x_jL_i.\end{gather}
由前两个恒等式易得
\begin{gather}
|E_+\phi(t,x)|\lesssim\frac1t|\Gamma\phi(t,x),\\
|E_-\phi(t,x)|\lesssim\frac1{|t-|x||}|\Gamma\phi(t,x)|,\end{gather}
其中 $|\Gamma\phi|=|S\phi|+|L\phi|$. 以及
\begin{equation}
|\bar\nabla\phi(t,x)|\lesssim\frac1t|\Gamma\phi(t,x)|.\end{equation}
显然还有
\begin{equation}
|\partial\phi(t,x)|\lesssim\frac1{|t-|x||}|\Gamma\phi(t,x)|,\end{equation}
或更一般地
\begin{equation}
|\partial^N\phi(t,x)|\lesssim\frac1{|t-|x||^N}|\Gamma^N\phi(t,x)|,\end{equation}
其中 $|\Gamma^N\phi|=\sum|\Gamma_1\ldots\Gamma_N\phi|$, $\Gamma_1,\ldots,\Gamma_N$ 为任意向量场 $S,L_1,\ldots,L_n$.
将上述不等式与我们所定义的范数结合, 推得
\begin{gather}
t\|E_+\phi(t)\|_{L^2}\lesssim E_1(\phi)^{\frac12},\\
t\|\nabla\phi(t)\|_{L^2}\lesssim E_1(\phi)^{\frac12},\\
\|uE_-\phi(t)\|_{L^2}\lesssim E_1(\phi)^{\frac12}.\end{gather}

现在我们准备证明如下结论:
\begin{proposition}\label{prop:4.15}
设 $\Box\phi=0$, 初数据满足上述假设. 则对所有 $t\geq\Lambda_0$, $s>\frac n2$,
\begin{gather}
\|\phi(t)\|_{L^\infty}\lesssim\Bigl(\frac1t\Bigr)^{\frac{n-1}{2}}E_s[\phi](t)^{\frac12},\\
\|(1+|u|)^k\partial^k\phi(t)\|_{L^\infty}\lesssim\Bigl(\frac1t\Bigr)^{\frac{n-1}{2}}E_{s+k}[\phi](t)^{\frac12}.\end{gather}
此外
\begin{gather}
\|E_+\phi(t)\|_{L^\infty}\lesssim\Bigl(\frac1t\Bigr)^{\frac{n+1}{2}}E_{s+1}[\phi](t)^{\frac12},\\
\|\bar\nabla\phi(t)\|_{L^\infty}\lesssim\Bigl(\frac1t\Bigr)^{\frac{n+1}{2}}E_{s+1}[\phi](t)^{\frac12},\\
\|(1+|u|)E_-\phi(t)\|_{L^\infty}\lesssim\Bigl(\frac1t\Bigr)^{\frac{n-1}{2}}E_{s+1}[\phi](t)^{\frac12}.\end{gather}
\end{proposition}

证明基于如下引理:
\begin{lemma}\label{lem:4.16}
设 $u(x)$ 为 $\mathbb{R}^n$($n\geq 2$)上光滑、具紧支集的函数. 则
\begin{equation}
|u(x)|\leq C\frac1{|x|^{n-1}}\Bigl(\|\partial_r u\|_{L^1}+\|(r\bar\nabla)^{n-1}\partial_r u\|_{L^1}\Bigr).\end{equation}
\end{lemma}

作为该引理的直接推论, 我们得到
\begin{equation}
|u(x)|\leq C\frac1{|x|^{n-1}}\sum_{k=0}^{n-1}\|\partial_r\Omega^k u\|_{L^1},\end{equation}
其中 $\Omega^k u$ 表示 $\Omega_1\ldots\Omega_k u$ 形式的所有导数, $\Omega_i=\epsilon_{ijk}x_j\partial_k$ 为角动量向量场. 现对固定某 $t$ 的函数 $u(x)=\phi^2(t,x)$ 应用不等式 (108):
\begin{equation}
|\phi(t,x)|^2\lesssim\frac1{|x|^{n-1}}\sum_{k=0}^{n-1}\sum_{i+j=k}\|\Omega^i\phi(t)\|_{L^2}\|\partial\Omega^j\phi(t)\|_{L^2}\lesssim\frac1{|x|^{n-1}}E_n[\phi].\end{equation}
由此推知, 当 $|x|\geq\frac t2$ 时,
\begin{equation}
|\phi(t,x)|\lesssim\Bigl(\frac1t\Bigr)^{\frac{n-1}{2}}E_n^{\frac12}[\phi].\end{equation}

现考虑区域 $|x|\leq\frac t2$. 取支在 $|x|\leq 1$ 内的检验函数 $\chi$, 且当 $|x|\leq\frac12$ 时 $\chi(x)=1$. 对函数 $u(x)=\chi(x/t)\phi^2(t,x)$ 应用 Sobolev 不等式
\begin{equation}
|u(x)|\lesssim\|D^n u\|_{L^1}.\end{equation}
于是当 $|x|\leq\frac t2$ 时, 顾及 (104) 再顾及 (31),
\begin{align*}
|\phi(t,x)|^2&\lesssim\sum_{i+j+k=n}\|D^i\chi(x/t)\|_{L^\infty}\cdot\|D^j\phi(t)\|_{L^2(|x|\leq\frac t2)}\cdot\|D^k\phi(t)\|_{L^2(|x|\leq\frac t2)}\\
&\lesssim\sum_{i+j+k=n}t^{-i}\cdot t^{-j}\|\Gamma^j\phi(t)\|_{L^2}\cdot t^{-k}\|\Gamma^k\phi(t)\|_{L^2}\\
&\lesssim t^{-n}E_n[\phi].\end{align*}
因此
\begin{equation}
\sup_{|x|\leq\frac t2}|\phi(t,x)|\lesssim\Bigl(\frac1t\Bigr)^{\frac n2}E_n^{\frac12}[\phi],\end{equation}
这与 (109) 一起证明了
\begin{equation}
\|\phi(t)\|_{L^\infty}\lesssim\Bigl(\frac1t\Bigr)^{\frac{n-1}{2}}E^{\frac12}[\phi],\end{equation}
正如所愿.
为证明\lemref{lem:4.16}, 在极坐标下写 $x=r\xi$, $\xi\in\mathbb{R}^{n-1}$:
\begin{equation}
u(r\xi)=-\int_r^\infty\partial_r u(\lambda\xi)\,d\lambda.\end{equation}
于是
\begin{equation}
\int_{|\xi|=1}|u(r\xi)|\,d\sigma(\xi)\leq\frac1{r^{n-1}}\int_{\mathbb{R}^n}|\partial_r u(y)|\,dy.\end{equation}
另一方面, 在单位球面 $S^{n-1}$ 上使用 Sobolev 不等式,
\begin{equation}
|u(x)|\leq c_n\Bigl(\|u(r\cdot)\|_{L^1(S^{n-1})}+\|(r\bar\nabla)^{n-1}u(r\cdot)\|_{L^1(S^{n-1})}\Bigr),\end{equation}
结合上面的不等式即证得所需结果.
\begin{definition}\label{def:4.17}
设 $\phi,\psi$ 为 $\mathbb{R}^{n+1}$ 上两个光滑函数. 我们定义
\begin{gather}
Q_0(\phi,\psi)=m^{\alpha\beta}\partial_\alpha\phi\partial_\beta\psi,\\
Q_{\alpha\beta}(\phi,\psi)=\partial_\alpha\phi\partial_\beta\psi-\partial_\alpha\psi\partial_\beta\phi.\end{gather}
\end{definition}

下面我们将利用前面建立的衰减估计, 并仿照 \cite{klai28} 的精神对零形式作适当分解, 来估计零二次型 $Q_0(\phi,\psi)$ 与 $Q_{ij}(\phi,\psi)$. 设 $\Box\phi=\Box\psi=0$. 考虑零形式 $Q_0(\phi,\psi)$ 并将其写成
\begin{equation}
Q_0(\phi,\psi)=-\frac12\bigl(E_+\phi\cdot E_-\psi+E_+\phi\cdot E_-\psi\bigr)+\bar\nabla\phi\cdot\bar\nabla\psi.\end{equation}
于是由\propref{prop:4.15} 的结果,
\begin{align*}
\|Q_0(\phi,\psi)(t)\|_{L^2}&\leq\frac12\Bigl(\|E_+(\phi)(t)\|_{L^\infty}\|E_-(\psi)(t)\|_{L^2}+\|E_-(\phi)(t)\|_{L^\infty}\|E_+(\psi)(t)\|_{L^2}\Bigr)\\
&\quad+\|\nabla\phi(t)\|_{L^\infty}\|\nabla\psi(t)\|_{L^2}\\
&\lesssim\Bigl(\frac1t\Bigr)^{\frac{n+1}{2}}E_{n+1}[\phi]^{\frac12}E_1[\psi]^{\frac12}.\end{align*}
类似地, 为估计 $Q_{ij}(\phi,\psi)$, 写 $\partial_i=A_i+\frac{x_i}{|x|}\partial_r$. 于是
\begin{equation}
Q_{ij}(\phi,\psi)=A_i\phi A_j\psi-A_j\phi A_i\psi+\partial_r\phi\Bigl(\frac{x_i}{r}A_j\phi-\frac{x_j}{r}A_i\Bigr)\phi.\end{equation}
顾及 (32) 中第三个公式以及 (31), 我们对 $Q_{ij}$ 得到同类型的估计.

\chapter{第二讲补编}

\section*{5.0.1 波映射}

从 Minkowski 时空到 Riemann 流形 $(M,g)$ 的波映射, 是泛函
\begin{equation}
\mathcal{L}[u]=\frac12\int_{\mathbb{R}^{1+n}}\langle du,du\rangle\,dt\,dx
\end{equation}
关于具紧支集变分的临界点, 其中在 $M$ 的局部坐标下 $\langle du,du\rangle=\sum_{\mu=0}^n\sum_{a,b}g_{ab}\partial_\mu u^a\partial^\mu u^b$. 该变分问题的 Euler--Lagrange 方程恰为 (WM) 的形式(在 $M$ 的局部坐标下), 其中 $\Gamma_{JK}^I$ 为 $M$ 在局部坐标系下的 Christoffel 符号, $N=\dim M$(参见例如 Shatah--Struwe \cite{klai58}).
\section*{5.0.2 Maxwell--Klein--Gordon 方程}

在下面的讨论中采用求和约定: 希腊指标从 $0$ 到 $n$ 求和, 罗马指标从 $1$ 到 $n$ 求和. 注意指标是相对 Minkowski 度量 $m_{\mu\nu}=\mathrm{diag}(-1,1,\ldots,1)$ 升降的. 例如 $\Box=\partial^\mu\partial_\mu$, $\Delta=\partial^j\partial_j$. 记 $\mathrm{i}$ 为虚数单位.
方程的未知量是定义于 Minkowski 时空上的一个一次形式 $A_\mu dx^\mu$(规范势)与一个标量 $\phi$:
\begin{equation}
\begin{aligned}
A_\mu&:\mathbb{R}^{1+n}\to\mathbb{R},\\
\phi&:\mathbb{R}^{1+n}\to\mathbf{C}.\end{aligned}
\end{equation}
电磁场为二次形式 $F_{\mu\nu}=\partial_\mu A_\nu-\partial_\nu A_\mu$. 相对规范势的协变导数为
\begin{equation}
D_\mu\phi=\partial_\mu\phi+\mathrm{i}A_\mu\phi.\end{equation}
我们寻找泛函
\begin{equation}
\mathcal{L}[A_\mu,\phi]=\int_{\mathbb{R}^{1+n}}\Bigl(-\frac14 F_{\mu\nu}F^{\mu\nu}-\frac12 D_\mu\phi\,\overline{D^\mu\phi}\Bigr)\,dt\,dx
\end{equation}
的临界点. 相应的 Euler--Lagrange 方程为
\begin{gather}
\partial^\mu F_{\mu\nu}=-\Im\bigl(\phi\,\overline{D_\nu\phi}\bigr),\\
D^\mu D_\mu\phi=0,\end{gather}
其中 $\Im z$ 表示 $z$ 的虚部.
设 $\chi$ 为 $\mathbb{R}^{1+n}$ 上实值函数, 考虑变换 $(A_\mu,\phi)\to(\widetilde A_\mu,\widetilde\phi)$, 由
\begin{gather}
\widetilde A_\mu=A_\mu-\partial_\mu\chi,\\
\widetilde\phi=e^{\mathrm{i}\chi}\phi
\end{gather}
给出. 显然, 电磁场在规范变换 $A_\mu\to\widetilde A_\mu$ 下保持不变; 且易见若 $(A_\mu,\phi)$ 满足 (MKG), 则 $(\widetilde A_\mu,\widetilde\phi)$ 亦满足(注意 $D_\mu$ 依赖于 $A_\mu$). 这就在满足 (MKG) 的数对 $(A_\mu,\phi)$ 的集合上给出了一个等价关系; 所谓 (MKG) 的一个解, 我们理解为此类数对的等价类.
于是我们拥有规范自由度; 即我们可自由选取给定解(等价类)的任一代表, 并可以规定规范势所应满足的条件. 传统的规范条件有:
\begin{itemize}
\item Lorentz 规范: $\partial^\mu A_\mu=0$;
\item Coulomb 规范: $\partial^j A_j=0$;
\item 时间规范: $A_0=0$.\end{itemize}

\paragraph*{(MKG) 的 Lorentz 规范.} 将 Lorentz 条件与 (MKG) 结合, 得到系统
\begin{gather}
\Box A_\mu=-\Im(\phi\,\overline{\partial_\mu\phi})+|\phi|^2A_\mu,\\
\Box\phi=-2\mathrm{i}A^\mu\partial_\mu\phi+A^\mu A_\mu\phi,\\
\partial^\mu A_\mu=0.\end{gather}
现注意, 若 $(A_\mu,\phi)$ 满足 (110a) 与 (110b), 其初数据
\begin{gather}
A_\mu|_{t=0}=a_\mu,\qquad \partial_t A_\mu|_{t=0}=b_\mu,\\
\phi|_{t=0}=\phi_0,\qquad \partial_t\phi|_{t=0}=\phi_1
\end{gather}
满足约束
\begin{equation}
b_0=\partial^j a_j,\qquad \Delta a_0-|\phi_0|^2 a_0=\partial^j b_j-\Im(\phi_0\overline{\phi}_1),\end{equation}
则 (110c) 自动满足. 因为由 (110a) 与 (110b), 量 $u=\partial^\mu A_\mu$ 满足 $\Box u=|\phi|^2u$, 且由 (111a) 与 (112) 有 $u|_{t=0}=\partial_t u|_{t=0}=0$. 由解的唯一性知 $u=0$.
因此, (110c) 等价于初数据上的约束 (112), 于是我们只需处理 (110a) 与 (110b). 从而 Lorentz 规范下的 (MKG) 在形式上即为 $\Box u=u\partial u+u^3$. 遗憾的是, 此类一般方程并不具备良好的局部正则性, 故 Lorentz 规范对我们的目的并不十分有用.
\paragraph*{(MKG) 的 Coulomb 规范.} 将 Coulomb 条件与 (MKG) 结合, 得到
\begin{gather}
\Delta A_0=-\Im(\phi\,\overline{\partial_t\phi})+|\phi|^2A_0,\\
\Box A_j=-\Im(\phi\,\overline{\partial_j\phi})+|\phi|^2A_j-\partial_j\partial_t A_0,\\
\Box\phi=-2\mathrm{i}A^j\partial_j\phi+2\mathrm{i}A_0\partial_t\phi+\mathrm{i}(\partial_t A_0)\phi+A^\mu A_\mu\phi,\\
\partial^j A_j=0.\end{gather}
此处我们将规范势分解为时间分量 $A_0$ 与空间分量 $A=A_j dx^j$. 在 $t=0$ 给定初数据:
\begin{gather}
A_j|_{t=0}=a_j,\qquad \partial_t A_j|_{t=0}=b_j,\\
\phi|_{t=0}=\phi_0,\qquad \partial_t\phi|_{t=0}=\phi_1.\end{gather}
对 $A_0$ 不加初值条件; 若记 $a_0=A_0|_{t=0}$, 则由 (113a) 有 $\Delta a_0-|\phi_0|^2a_0=-\Im(\phi_0\overline{\phi}_1)$.
若数据散度自由, 即
\begin{equation}
\partial^j a_j=\partial^j b_j=0,\end{equation}
则 (113d) 自动满足. 因为若 $(A_0,A,\phi)$ 满足 (113a)--(113c), 则 $u=\partial^j A_j$ 满足 $\Box u=|\phi|^2u$, 而若 (114) 与 (115) 成立, 则 $u|_{t=0}=\partial_t u|_{t=0}=0$.
于是我们只需方程 (113a)--(113c). 其中第一个为椭圆方程, 较易处理, 故将其排除在我们的模型方程之外. 剩余两个方程右端含有三类项:
\begin{itemize}
\item 涉及 $A_0$ 的"椭圆项", 统称 $\mathcal{E}$;
\item $A_j$ 与 $\phi$ 中的三次项, 统称 $\mathcal{C}$;
\item 具零形式结构的二次项.\end{itemize}

归入后一类的项为 $-\Im(\phi\,\overline{\partial_j\phi})$ 与 $-2\mathrm{i}A^j\partial_j\phi$. 我们现在揭示这些表达式内蕴的零形式结构(源于 Coulomb 条件).
将 $\phi$ 分解为实部与虚部: $\phi=u+\mathrm{i}v$. 则
\begin{equation}
-\Im(\phi\,\overline{\partial_j\phi})=u\partial_j v-v\partial_j u,\end{equation}
于是 (113b) 作为(依赖于时间的)一次形式方程在 $\mathbb{R}^n$ 上可写为
\begin{equation}
\Box A=u\,dv-v\,du+\mathcal{C}+\mathcal{E}.\end{equation}
两边同时作用外微分 $d$:
\begin{equation}
\Box(dA)=2\,du\wedge dv+d\mathcal{C}.\end{equation}
而
\begin{equation}
du\wedge dv=\frac12 Q_{jk}(u,v)\,dx^j\wedge dx^k,\end{equation}
由此
\begin{equation}
\Box F_{jk}=Q_{jk}(u,v)+\partial\mathcal{C}.\end{equation}
Coulomb 规范条件推出 $\partial^k F_{jk}=-\Delta A_j$, 于是
\begin{equation}
-\Delta\Box A_j=\partial^k Q_{jk}(u,v)+\partial^2\mathcal{C}.\end{equation}
从而在 Riesz 算子意义下,
\begin{equation}
\Box A=D^{-1}Q(\Re\phi,\Im\phi)+\mathcal{C},\end{equation}
其中 $Q$ 为零形式 $Q_{jk}$ 的某个线性组合. 由于三次项 $\mathcal{C}$ 更易估计, 我们将其排除在模型问题之外.
现考虑方程 (113c). 分离实部与虚部, 有
\begin{equation}
\begin{aligned}
\Box u&=2A\cdot\nabla v+\mathcal{C}+\mathcal{E},\\
\Box v&=-2A\cdot\nabla u+\mathcal{C}+\mathcal{E}.\end{aligned}
\end{equation}
(此处我们将 $A$ 借指标提升视为向量场; $\nabla$ 表示空间变量的梯度. )我们断言, 项 $A\cdot\nabla u$ 与 $A\cdot\nabla v$ 具有零形式结构, 这是因为 $A$ 无散度(由 Coulomb 条件). 设 $B_{jk}$ 为方程
\begin{equation}
\Delta B_{jk}=\partial_j A_k-\partial_k A_j
\end{equation}
的唯一解(在适当的正则性假设下). 由 Coulomb 条件,
\begin{equation}
\Delta\partial^j B_{jk}=\Delta A_k,\qquad\text{即}\quad \partial^j B_{jk}=A_k.\end{equation}
于是
\begin{equation}
A\cdot\nabla u=\partial^j B_{jk}\partial^k u=\frac12 Q_{jk}(u,B^{jk}).\end{equation}
上述关于 $u=\Re\phi$ 与 $v=\Im\phi$ 的方程因而可重写为
\begin{equation}
\begin{aligned}
\Box\Re\phi&=Q_{jk}(\Im\phi,B^{jk})+\mathcal{C}+\mathcal{E},\\
\Box\Im\phi&=Q_{jk}(B^{jk},\Re\phi)+\mathcal{C}+\mathcal{E}.\end{aligned}
\end{equation}
但由 (117), $B$ 在 Riesz 算子意义下形如 $D^{-1}A$. 将此与 (116) 结合, 并处处舍去项 $\mathcal{C}$ 与 $\mathcal{E}$, 我们得到形如 ("MKG") 的一个系统, 此即我们的 (MKG) 模型.
\section*{5.0.3 Yang--Mills 方程}

设 $G$ 为经典紧 Lie 矩阵群之一(如 $SO(k,\mathbb{R})$ 或 $SU(k,\mathbf{C})$), $\mathfrak{g}$ 为其 Lie 代数. 未知量为 $\mathbb{R}^{1+n}$ 上取值于 $\mathfrak{g}$ 的一次形式 $A_\mu dx^\mu$. 相应的协变导数为
\begin{equation}
D_\mu H=\partial_\mu H+[A_\mu,H],\end{equation}
其中 $H$ 为任意取值于 $\mathfrak{g}$ 的张量场, $[\cdot,\cdot]$ 为矩阵交换子. 曲率(曲率为取值于 $\mathfrak{g}$ 的二次形式)为
\begin{equation}
F_{\mu\nu}=\partial_\mu A_\nu-\partial_\nu A_\mu+[A_\mu,A_\nu].\end{equation}
作用量为
\begin{equation}
\mathcal{L}[A_\mu]=-\frac14\int\langle F_{\mu\nu},F^{\mu\nu}\rangle\,dt\,dx,\end{equation}
其中 $\langle\cdot,\cdot\rangle$ 为 $\mathfrak{g}$ 上继承自环境空间的内积(例如 $SO(k,\mathbb{R})$ 嵌入 $\mathbb{R}^{k^2}$, 其 Lie 代数可视为后者的子空间). Euler--Lagrange 方程为
\begin{equation}
D^\nu F_{\mu\nu}=0.\end{equation}
设 $O$ 为 $\mathbb{R}^{1+n}$ 上取值于 $G$ 的函数. 考虑规范变换 $A_\mu\to\widetilde A_\mu$, 由
\begin{equation}
\widetilde A_\mu=O A_\mu O^{-1}-\partial_\mu O\,O^{-1}
\end{equation}
给出. 计算表明曲率随之变换为
\begin{equation}
\widetilde F_{\mu\nu}=O F_{\mu\nu}O^{-1}.\end{equation}
记 $\widetilde D_\mu$ 为对应于 $\widetilde A_\mu$ 的协变导数, 则
\begin{equation}
\widetilde D^\nu\widetilde F_{\mu\nu}=O D^\nu F_{\mu\nu}O^{-1},\end{equation}
故 (119) 在规范变换下不变. 于是我们拥有规范自由度, 并可对 $A_\mu$ 施加规范条件.
\paragraph*{(YM) 的 Coulomb 规范.} 相对 Coulomb 条件 $\partial^j A_j=0$, (119) 取形式(见 \cite{klai34})
\begin{gather}
\Delta A_0=2[\partial^j A_0,A_j]+[A^j,\partial_t A_j]+[A^j,[A_0,A_j]],\\
\Box A_j+\partial_t\partial_j A_0=-2[A^k,\partial_k A_j]+[A^k,\partial_j A_k]+[\partial_t A_0,A_j]+2[A_0,\partial_t A_j]\nonumber\\
\qquad-\,[A_0,\partial_j A_0]-[A^k,[A_k,A_j]]+[A_0,[A_0,A_j]],\\
\partial^j A_j=0.\end{gather}
遗憾的是, 假设存在整体 Coulomb 规范, 会对初数据施加一个限制性的小性假设. 在 \cite{klai34} 中, 这一困难借助局部论证得以解决. 依 \cite{klai47}, 我们忽略这一复杂之处, 而从系统 (119) 导出我们的模型方程.
与 (MKG) 的讨论一样, (119c) 化为初数据上的一个约束. 关于 $A_0$ 的方程是椭圆的, 故我们忽略它. 至于 (119b), 我们只保留右端的前两项, 因为其余项要么涉及 $A_0$(对其我们期望得到比 $A_j$ 更好的估计), 要么为三次项.
现将 (119b) 写为(依赖于时间的、取值于 $\mathfrak{g}$ 的)$\mathbb{R}^n$ 上一次形式的方程(忽略右端除前两项外的一切):
\begin{equation}
\Box A+d(\partial_t A_0)=S+T,\end{equation}
其中 $A=A_j dx^j$, $S=-2[A^k,\partial_k A_j]dx^j$, $T=[A^k,\partial_j A_k]dx^j$. 两边同时作用外微分 $d$:
\begin{equation}
\Box dA=dS+dT.\end{equation}
设 $B$ 为 (117) 所确定的二次形式(此处取值于 $\mathfrak{g}$). 于是
\begin{equation}
\Box\Delta B_{jk}=\partial_j S_k-\partial_k S_j+\partial_j T_k-\partial_k T_j.\end{equation}
由 (118) 推得
\begin{equation}
-\Box\Delta A_j=\partial^k(\partial_j S_k-\partial_k S_j+\partial_j T_k-\partial_k T_j),\end{equation}
故就频率空间中的估计而言, 我们可用下式替代 (119b):
\begin{equation}
\Box A=S+D^{-1}dT.\end{equation}
剩下要辨明隐藏于 $S$ 与 $dT$ 中的零形式结构. 首先,
\begin{align*}
S_j&=[\partial_k B^{kl},\partial_l A_j]=\frac12[\partial_k B^{kl}-\partial_k B^{lk},\partial_l A_j]\\
&=\frac12\bigl([\partial_k B^{kl},\partial_l A_j]-[\partial_l B^{kl},\partial_k A_j]\bigr)\\
&=\frac12\bigl(\partial_k B^{kl}\partial_l A_j-\partial_l B^{kl}\partial_k A_j+\partial_k A_j\partial_l B^{kl}-\partial_l A_j\partial_k B^{kl}\bigr),\end{align*}
故矩阵 $S_j$ 的每一项都是形如 $Q_{kl}(A,B)$ 的项的线性组合, 其中 $A$ 与 $B$ 表示 $A_j$ 与 $B_{kl}$ 的任意两个分量. 但由 (118), 我们可用 $D^{-1}A$ 替代 $B$. 概略地,
\begin{equation}
S\sim Q(A,D^{-1}A).\end{equation}
现考虑一次形式 $T$. 我们计算:
\begin{align*}
(dT)_{jk}&=\partial_j[A^l,\partial_k A_l]-\partial_k[A^l,\partial_j A_l]\\
&=[\partial_j A^l,\partial_k A_l]-[\partial_k A^l,\partial_j A_l].\end{align*}
于是矩阵 $(dT)_{jk}$ 的每一项都是形如 $Q_{jk}(A,A')$ 的项的线性组合, 其中 $A$ 与 $A'$ 表示 $A_l$($1\leq l\leq n$)的任意两个分量. 将此与 (121) 及 (120) 结合, 我们即得 Yang--Mills 方程的模型 ("YM").

\chapter{第三讲补编}

此结果的一个重要推论如下:
\begin{theorem}\label{thm:6.1}
对任意渐近平直初数据集 $(\Sigma,\bar g,\bar k)$, 其极大未来发展 $(\mathcal{M},g)$ 必可找到一个具有紧闭包的区域 $\Omega_0\subset\Sigma$, 使得其在 $\mathcal{M}$ 中影响域 $\mathcal{C}^+(\Omega_0)$ 的边界 $\mathcal{D}_l^+$ 拥有完备的零测地线生成元.\end{theorem}

上述推论可用来引入"完备未来零无穷远"的概念:
\begin{definition}\label{def:6.2}
若对任一正实数 $A$, 我们都能找到一个包含定理 1.2 中区域 $\Omega_0$ 的区域 $\Omega$, 使得 $\Omega$ 在 $\mathcal{M}$ 中依赖域的边界 $\mathcal{D}^-$具有如下性质: 它的每条零测地线生成元, 从 $\mathcal{D}^-\cap\mathcal{D}_0^+$ 起算的总仿射长度至少为 $A$, 则称渐近平直初数据集 $(\Sigma,\bar g,\bar k)$ 的极大未来发展 $(\mathcal{M},g)$ 拥有\textbf{完备未来零无穷远}.\end{definition}

对于充分大的初数据集, 奇点的不可避免出现, 促使 Penrose 提出两个猜想, 即弱宇宙监督猜想与强宇宙监督猜想.
\noindent\textbf{猜想 1(弱宇宙监督猜想 WCC)}(\textit{Conjecture 1, Weak Cosmic Censorship})．\label{conj:wcc}一般的渐近平直初数据, 其极大未来发展拥有完备的未来零无穷远.
\par\medskip

WCC 猜想并不排除奇点可能被局部观测者看到的可能性. 这可能导致如下悖论式的情形: 此类观测者所作观测的结果缺乏唯一可预测性. 强宇宙监督猜想正是为了禁止奇点这一不良特性而提出的.
\noindent\textbf{猜想 2(强宇宙监督猜想 SCC)}(\textit{Conjecture 2, Strong Cosmic Censorship})．\label{conj:scc}一般的初数据集, 其极大未来发展作为 Lorentz 流形, 在连续意义下局部不可延拓.
\par\medskip

更技术地说, 这意味着除去某些可能的例外初值条件外, 一个初数据集的极大未来发展具有如下性质: 沿任一未来、不可延拓、有限长度的类时测地线, 时空曲率分量(相对沿该测地线平行移动的标准正交标架表示)当弧长趋近于其极限值时必趋于无穷.
\section*{6.3 \thmref{thm:3.10} 到\thmref{thm:3.12} 的化简}

下面我们描述\thmref{thm:3.10}(以下记为定理 (A))到\thmref{thm:3.12} 之衰减估计(以下记为定理 (B))证明的化简. 前五个步骤如今已是标准的, 参见 \cite{klai2}, 尤其是 \cite{klai72}; 最后一步的化简则是 \cite{klai29} 几何方法的典型特征.
\section*{6.4 步骤 1: 能量估计}

如前所述, 我们对拟线性波动方程的能量估计表明: 只要 $\|\partial\phi\|_{L_{[0,T]}^1L_x^\infty}\leq 1$, 解 $\phi$ 的 $L_{[0,T]}^\infty H^s$ 范数即被初数据 $\phi[0]$ 的 $H^s$ 范数控制. 下面的命题是这一陈述的更精确版本.
\setcounter{mtthm}{4}% 6.3、6.4 为步骤标题, 下一定理环境显示 6.5

\begin{proposition}[能量估计]\label{prop:6.5}
设 $\phi\in C([0,T],H^s)\cap C^1([0,T],H^{s-1})$ 为区间 $[0,T]$ 上 (67) 的解, 对某 $s\geq 1$ 满足 $\|\phi\|_{L_{[0,T]}^\infty L_x^\infty}\leq\Lambda_0$. 则 $\phi$ 满足如下能量估计:
\begin{equation}
\|\phi\|_{L_{[0,T]}^\infty\dot H^s}\leq C\bigl(\|\partial\phi\|_{L_{[0,T]}^1L_x^\infty},\Lambda_0\bigr)\|\phi[0]\|_{\dot H^s}.\end{equation}
此处常数 $C(a,b)$ 表示对 $a,b$ 的依赖, $\dot H^s$ 为通常的齐次 Sobolev 空间.\end{proposition}

\section*{6.6 步骤 2: 化归 Strichartz 型估计}

由 Cauchy--Schwarz 不等式, $\|\partial\phi\|_{L_{[0,T]}^1L_x^\infty}\leq T^{\frac12}\|\partial\phi\|_{L_{[0,T]}^2L_x^\infty}$. 于是 Strichartz 不等式 (71) 与时间区间 $[0,T]$ 之小性, 可用来在 $H^{s_*}$ 空间中闭合能量估计. 这就有效地给出了该问题所需的局部适定性. 定理 (A) 从而可化归如下自举论证.
\setcounter{mtthm}{6}% 6.6 为步骤标题, 下一定理环境显示 6.7

\begin{theorem}[A1]\label{thm:6.7}
设 $\phi\in C([0,T],H^{s_*})\cap C^1([0,T],H^{s_*-1})$ 为区间 $[0,T]$($T\leq 1$)上 (67) 的解. 假定
\begin{equation}
\|\phi\|_{L_{[0,T]}^\infty H^{s_*}}+\|\partial\phi\|_{L_{[0,T]}^2L_x^\infty}\leq B_0,\end{equation}
其中常数 $B_0\leq c_{s_*}^{-1}\Lambda_0$, $c_{s_*}$ 为嵌入 $H^{s_*}\subset L^\infty$ 的 Sobolev 常数. 则 $\phi$ 满足如下时间局部的 Strichartz 型估计:
\begin{equation}
\|\partial\phi\|_{L_{[0,T]}^2L_x^\infty}\leq C(B_0)T^{s_*-s_0}\|\phi\|_{L_{[0,T]}^\infty H^{s_*}}.\end{equation}
\end{theorem}

\setcounter{mtthm}{8}% 6.8 为步骤标题, 下一定理环境显示 6.9

\section*{6.8 步骤 3: Strichartz 型估计的二进版本与仿微分逼近}

为证明 Strichartz 型估计 (124), 我们可将方程 (67) 视为关于 $\phi$ 的、具粗糙系数的线性波动方程. 对系数作磨光, 并研究与一参数相关的、具光滑系数的线性波动方程族, 是有利的. 首先引入 $\phi^\lambda$: 将 $\phi$ 限制于傅里叶空间中频率 $\sim\lambda$ 的二进块得到. 更确切地, 设 $\zeta$ 为支在球壳 $\{\xi:\frac12\leq|\xi|\leq 2\}$ 上的光滑函数(此处 $\xi$ 表示空间傅里叶变换的变量), 并满足 $\sum_{k\in\mathbb{Z}}\zeta(2^k\xi)=1$($\forall\xi\in\mathbb{R}^n\setminus\{0\}$). 令 $\lambda$ 为二进参数 $\lambda=2^k$, $k\in\mathbb{Z}$, 记 $P_\lambda$ 为"投影子":
\begin{equation}
P_\lambda f(x)=f^\lambda(x)=\int e^{-\mathrm{i}x\cdot\xi}\zeta(\lambda^{-1}\xi)\hat f(\xi)\,d\xi.\end{equation}
又定义
\begin{equation}
f_{\leq\lambda}=S_\lambda f=\sum_{\mu\leq\lambda}P_\mu f.\end{equation}

定理 (A1) 由如下关于 $\phi^\lambda=P_\lambda\phi$ 的 Strichartz 型估计的二进版本推出.
\begin{theorem}[A2]\label{thm:6.9}
设 $\phi$ 如定理 (A1). 固定一大参数 $\Lambda$. 则对每个 $\lambda\geq\Lambda$, $\phi^\lambda$ 满足 Strichartz 型估计
\begin{equation}
\|\partial\phi^\lambda\|_{L_{[0,T]}^2L_x^\infty}\leq C(B_0)c_\lambda T^{s_*-s_0}\|\phi\|_{L_{[0,T]}^\infty H^{s_*}},\end{equation}
其中常数 $c_\lambda$ 满足 $\sum_\lambda c_\lambda^2\leq 1$. 类似估计对 $\phi_{\leq\Lambda}$ 亦成立.\end{theorem}

\textbf{注 (A2).} 对低频部分, 估计 (125) 对 $\phi_{\leq\Lambda}$ 由 Sobolev 不等式显然成立:
\begin{equation}
\|\partial\phi_{\leq\Lambda}\|_{L_{[0,T]}^2L_x^\infty}\leq cT^{\frac12}\|\phi_{\leq\Lambda}\|_{L_{[0,T]}^\infty H^{\frac52+\epsilon}}\leq c\Lambda^{\frac52+\epsilon-s_*}T^{\frac12}\|\phi\|_{L_{[0,T]}^\infty H^{s_*}},\end{equation}
其中 $c$ 为嵌入 $H^{\frac32+\epsilon}(\mathbb{R}^3)\subset L^\infty(\mathbb{R}^3)$ 的范数. 由于 $s_*$ 假定充分接近 $s_0=2+\frac12(2-\sqrt3)$, 且 $\Lambda$ 为可仅依赖于 $B_0$ 的固定大参数, 我们便得到 $\phi$ 低频部分所需的界.
我们只关注大频率 $\lambda\geq\Lambda$. 下一命题表明, 每个 $\phi^\lambda$ 满足一个具光滑度量 $g_{\leq\lambda^a}^{ij}=S_{\lambda^a}g^{ij}(S_\lambda^a\phi)$ 的非齐次波动方程, 其中参数 $a\in[0,1]$ 取固定值.
\begin{proposition}\label{prop:6.10}
设 $\phi$ 如定理 (A1). 固定参数 $a$ 的值, $a\in[0,1]$. 则对每个 $\lambda\geq\Lambda$, $\phi^\lambda$ 满足方程
\begin{equation}
\begin{aligned}
\Box_{g_{\leq\lambda^a}}\phi^\lambda&=-\partial_t^2\phi^\lambda+g_{\leq\lambda^a}^{ij}\partial_i\partial_j\phi^\lambda=R_\lambda^a,\\
\phi^\lambda|_{t=0}&=\phi_0^\lambda,\qquad \partial_t\phi^\lambda|_{t=0}=\phi_1^\lambda.\end{aligned}
\end{equation}
此外, 右端 $R_\lambda^a$ 的傅里叶支集含于集合 $\{\xi:\lambda\leq|\xi|\leq 4\lambda\}$ 中, 且对所有实数 $s>1$ 与任意 $t\in[0,T]$,
\begin{gather}
\|R_\lambda(t)\|_{\dot H^s}\leq c\lambda^s\|R_\lambda\|_{L_x^2}\lesssim C(B_0)\lambda^{1-a}c_\lambda\|\partial\phi\|_{L_x^\infty}\|\phi\|_{\dot H^s},\\
c_\lambda:\ \sum_\lambda c_\lambda^2\leq 1.\end{gather}
\end{proposition}

\setcounter{mtthm}{11}% 6.11 为步骤标题, 下一定理环境显示 6.12

\section*{6.11 步骤 4: 频率相关区间上的 Strichartz 估计}

此步骤将定理 (A) 的证明化归为如下线性化方程 $\Box_{g_{\lambda^a}}\psi=0$ 在小时间区间 $I$(大小 $\approx T\lambda^{-(1-a)}$)上精确 Strichartz 估计的证明. 最终 Strichartz 估计中正则性的损失, 正是将 $[0,T]$ 上这些尖锐的 Strichartz 估计对区间 $I$ 求和的结果.
\begin{theorem}[A3]\label{thm:6.12}
设 $\lambda\geq\Lambda$, $\psi$ 为线性波动方程 $\Box_{g_{\lambda^a}}\psi=0$ 的解, 初数据 $\psi[0]$ 满足 $\operatorname{supp}\hat\psi[0]\in\{\frac12\lambda\leq|\xi|\leq 2\lambda\}$, 度量 $g_{\lambda^a}$ 如\propref{prop:6.10} 所定义. 固定参数 $a$ 的值, $a=\frac{5-2s_0}{1-2s_0+2s_*}$ 且 $a<-1+\sqrt3$. 则存在区间 $[0,T]$ 的一个分拆 $\{I\}$ 为子区间 $I$, 使得每个 $I$ 的大小 $|I|\leq T\lambda^{-(1-a)}$, 子区间的个数约为 $\lambda^{1-a}$, 且在每个 $I$ 上, $\psi$ 满足一带任意固定小量 $\epsilon>0$ 的 Strichartz 估计:
\begin{equation}
\|P_\lambda\partial\psi\|_{L_I^2L_x^\infty}\leq C(B_0)|I|^\epsilon\|\psi[0]\|_{\dot H^{2+\epsilon}}.\end{equation}
\end{theorem}

\textbf{注 (A3).} 依自举假设 (123), 拟线性问题的解 $\phi$ 满足估计 $\|\partial\phi\|_{L_{[0,T]}^2L_x^\infty}\leq B_0$. 由此易见存在时间区间 $[0,T]$ 的一个子分拆 $\{I\}$, 其总子区间数介于 $\lambda^{1-a}$ 与 $2\lambda^{1-a}$ 之间, 每个 $I$ 的大小不超过 $T\lambda^{-(1-a)}$, 且在每个 $I$ 上我们有
\begin{equation}
\|\partial\phi\|_{L_I^2L_x^\infty}\leq\lambda^{-\frac{1-a}{2}}\|\partial\phi\|_{L_{[0,T]}^2L_x^\infty}.\end{equation}
这一构造定义了定理 (A3) 中提到的子分拆 $\{I\}$.
\section*{6.13 度量 $g_{\lambda^a}$ 的性质}

不等式 (128) 是具变系数(度量)波动方程 $\Box_{g_{\leq\lambda^a}}\psi=0$ 在频率相关区间 $I$ 上的 Strichartz 估计.
度量 $g_{\leq\lambda^a}=S_{\lambda^a}g(S_{\lambda^a}\phi)$ 依赖于拟线性问题的解 $\phi$. 下一命题陈述了由 $\phi$ 上的自举条件 (123) 与注 (A3) 中区间分拆的构造所导出的族 $g_{\leq\lambda^a}$ 的性质.
\setcounter{mtthm}{13}% 6.13 为步骤标题, 下一定理环境显示 6.14

\begin{proposition}\label{prop:6.14}
设 $\phi\in C([0,T],H^{s_*})\cap C^1([0,T],H^{s_*-1})$ 为区间 $[0,T]$($T\leq 1$)上 (67) 的解. 假设 $\phi$ 满足定理 (A1) 的假设 (123). 考虑由注 (A3) 定义的子分拆 $\{I\}$. 则度量族 $g_{\leq\lambda^a}=S_{\lambda^a}g(S_{\lambda^a}\phi)$ 满足如下条件:
对所有子区间 $I$ 与所有非负整数 $m$,
\begin{gather}
\|\partial^{1+m}g_{\leq\lambda^a}\|_{L_I^1L_x^\infty}\leq\lambda^{-(1-a)+am}\bar B_0,\\
\|\partial^{1+m}g_{\leq\lambda^a}\|_{L_I^2L_x^\infty}\leq\lambda^{-\frac{1-a}{2}+am}\bar B_0,\\
\|\partial^{1+m}g_{\leq\lambda^a}\|_{L_I^\infty L_x^\infty}\leq\lambda^{\frac{a^2}{2}+am}\bar B_0,\\
\|\partial_x^{\frac12+m}(\partial^2 g_{\leq\lambda^a})\|_{L_I^\infty L_x^2}\leq\lambda^{\frac{a^2}{2}+am}\bar B_0,\\
\|\partial^m\Box_{g_{\leq\lambda^a}}g_{\leq\lambda^a}\|_{L_I^1L_{\bar x}^\infty}\leq\lambda^{-(1-a)+am}\bar B_0.\end{gather}
常数 $\bar B_0$ 只依赖于常数 $M_0$ 与 $B_0$.\end{proposition}

注: 由构造可知, 度量 $g_{\leq\lambda^a}$ 的频率只在相对空间变量 $x$ 的傅里叶对偶变量方向被截断于 $\lambda^a$ 之上. 因此, 关于 $x$ 的每一次微分至多引入一个因子 $\lambda^a$. 然而, 利用 $g_{\leq\lambda^a}$ 依赖于波动方程之解这一事实, 对时间导数我们也能得到同样的结论.
定理 (A3) 可重新表述为: 关于具背景度量 $g_{\leq\lambda^a}$(满足\propref{prop:6.14} 之估计)的线性波动方程, 在固定小区间 $I$ 上局部 Strichartz 估计的结果.
\begin{theorem}[A4]\label{thm:6.15}
设 $\psi$ 为线性波动方程
\begin{equation}
\begin{aligned}
\Box_{g_{\leq\lambda^a}}\psi&=-\partial_t^2\psi+g_{\leq\lambda^a}^{ij}\partial_i\partial_j\psi=0,\\
\psi|_{t=0}&=\psi_0,\qquad \partial_t\psi|_{t=0}=\psi_1
\end{aligned}
\end{equation}
在区间 $I$(长度 $|I|\leq\lambda^{-(1-a)}$)上的解, 初数据 $\psi[0]$ 在傅里叶空间中支于集合 $\{\xi:\frac12\lambda\leq|\xi|\leq 2\lambda\}$. 假设度量 $g_{\lambda^a}$ 满足\propref{prop:6.14} 的 (130)--(134), 参数 $a$ 取 $a<-1+\sqrt3$. 令 $P_\lambda$ 为到傅里叶空间集合 $\{\xi:\frac12\lambda\leq|\xi|\leq 2\lambda\}$ 的投影. 则对充分大参数 $\Lambda$、所有二进 $\lambda\geq\Lambda$ 与固定 $\epsilon>0$,
\begin{equation}
\|P_\lambda\partial\psi\|_{L_I^2L_x^\infty}\leq C(\bar B_0)|I|^\epsilon\|\psi[0]\|_{\dot H^{2+\epsilon}},\end{equation}
其中常数 $C(\bar B_0)$ 与 $\lambda$ 无关.\end{theorem}

\setcounter{mtthm}{16}% 6.16 为步骤标题, 下一定理环境显示 6.17

\section*{6.16 步骤 5: 重标度}

将问题 (135) 替换为它的重标度版本是方便的, 使得初数据的频率支集满足 $|\xi|\sim 1$, 且时间区间 $I$ 的长度为 $\leq\lambda^a$.
必要时对问题作时间平移, 可设时间区间 $I$ 始于 $t=0$. 引入重标度度量族 $h_\lambda$:
\begin{equation}
h_\lambda(t,x)=g_{\leq\lambda^a}(\lambda^{-1}t,\lambda^{-1}x).\end{equation}
\propref{prop:6.14} 推出, $h_\lambda$ 在时间区间 $I=[0,t_*]$($t_*\leq\lambda^a$)上满足如下估计:
\begin{gather}
\|\partial^{1+m}h_\lambda\|_{L_{I_\lambda}^1L_x^\infty}\lesssim\lambda^{-(1-a)(m+1)},\\
\|\partial^{1+m}h_\lambda\|_{L_{I_\lambda}^2L_x^\infty}\lesssim\lambda^{-\frac{2-a}{2}-(1-a)m},\\
\|\partial^{1+m}h_\lambda\|_{L_{I_\lambda}^\infty L_x^\infty}\lesssim\lambda^{-1+\frac{a^2}{2}-(1-a)m},\\
\|\partial_x^{\frac12+m}(\partial^2 h_\lambda)\|_{L_{I_\lambda}^\infty L_x^2}\lesssim\lambda^{-1+\frac{a^2}{2}-(1-a)m},\\
\|\partial^m\Box_{h_\lambda}h_\lambda\|_{L_{I_\lambda}^1L_x^\infty}\lesssim\lambda^{-(2-a)-(1-a)m}.\end{gather}

重标度后, 定理 (A4) 化为

\begin{theorem}[A5]\label{thm:6.17}
设 $\psi$ 为线性波动方程
\begin{equation}
\begin{aligned}
\Box_{h_\lambda}\psi&=-\partial_t^2\psi+h_\lambda^{ij}\partial_i\partial_j\psi=0,\\
\psi|_{t=0}&=\psi_0,\qquad \partial_t\psi|_{t=0}=\psi_1
\end{aligned}
\end{equation}
在区间 $[0,t_*]$($t_*\leq\lambda^a$)上的解. 假设参数 $\lambda\geq\Lambda$($\Lambda$ 为充分大常数), 且度量 $h_\lambda$ 满足 (138)--(142), 参数 $a$ 满足 $a<-1+\sqrt3$. 令 $P$ 为到傅里叶空间集合 $\{\xi:1\leq|\xi|\leq 2\}$ 的投影算子. 则
\begin{equation}
\|P\partial\psi\|_{L_{[0,t_*]}^2L_x^\infty}\lesssim|t_*|^\epsilon\bigl(\|\partial\psi_0\|_{L_x^2}+\|\psi_1\|_{L_x^2}\bigr),\end{equation}
其中不等式 $\lesssim$ 中的常数与 $\lambda$ 无关.\end{theorem}

注: 定理 (A5) 不含关于初数据 $\psi[0]$ 傅里叶支集的任何假设.
\section*{6.18 步骤 6: 衰减估计}

标准 $TT^*$ 型论证的一个变体(见 \cite{klai29})使我们能将 Strichartz 估计 (144) 化归为相应的色散不等式, 即 (145). 在此过程中, 我们将方程 $\Box_{h_\lambda}\psi=0$ 替换为几何波动方程
\begin{equation}
\Box_{h_\lambda}\psi=-\frac{1}{\sqrt{\det h_\lambda}}\partial_t\sqrt{\det h_\lambda}\,\partial_t\psi+\frac{1}{\sqrt{\det h_\lambda}}\partial_i\bigl(h_\lambda^{ij}\sqrt{\det h_\lambda}\,\partial_j\psi\bigr)=0.\end{equation}

\setcounter{mtthm}{18}% 6.18 为步骤标题, 下一定理环境显示 6.19

\begin{theorem}[A6]\label{thm:6.19}
设 $\psi$ 为线性波动方程
\begin{equation}
\begin{aligned}
\Box_{h_\lambda}\psi&=-\frac{1}{\sqrt{\det h_\lambda}}\partial_t\sqrt{\det h_\lambda}\,\partial_t\psi+\frac{1}{\sqrt{\det h_\lambda}}\partial_i\bigl(h^{ij}_\lambda\sqrt{\det h_\lambda}\,\partial_j\psi\bigr)=0,\\
\psi|_{t=0}&=\psi_0,\qquad \partial_t\psi|_{t=0}=\psi_1
\end{aligned}
\end{equation}
在区间 $[0,t_*]$($t_*\leq\lambda^a$)上的解, 初数据 $\psi[0]$ 在傅里叶空间中支于集合 $\{\xi:\frac12\leq|\xi|\leq 2\}$. 我们只考虑大参数值 $\lambda\geq\Lambda$. 假设度量 $h_\lambda$ 满足 (138)--(142), 参数 $a$ 满足 $a<-1+\sqrt3$. 则对所有 $t\leq t_*$ 与任意固定小量 $\epsilon>0$,
\begin{equation}
\|P\partial\psi(t)\|_{L_x^\infty}\lesssim\frac{1}{(1+|t|)^{1-\epsilon}}\|\psi[0]\|_{L_x^1}.\end{equation}
\end{theorem}

我们通过在物理空间将初数据 $\psi[0]$ 分解为支集本质不交、含于半径 $\frac12$ 的球中的函数之和, 作最后的化简. 利用 $L^1$ 范数的可加性与标准 Sobolev 不等式, 我们可将色散不等式 (146) 化归为 $L^2$--$L^\infty$ 衰减估计.
\begin{theorem}[B]\label{thm:6.20}
设 $\psi$ 为线性波动方程 (145) 在区间 $[0,t_*]$($t_*\leq\lambda^a$)上的解, 初数据 $\psi[0]$ 支于物理空间中以原点为中心、半径 $\frac12$ 的球 $B_{\frac12}(0)$ 内. 我们固定一大常数 $\Lambda$, 并只考虑大参数值 $\lambda\geq\Lambda$. 假设度量 $h_\lambda$ 满足 (138)--(142), 参数 $a$ 满足 $a<-1+\sqrt3$. 则对所有 $t\leq t_*$、任意小量 $\epsilon>0$ 与充分大整数 $m>0$,
\begin{equation}
\|P\partial\psi(t)\|_{L_x^\infty}\lesssim\frac{1}{(1+|t|)^{1-\epsilon}}\sum_{k=1}^m\|\partial^k\psi[0]\|_{L_x^2}.\end{equation}
\end{theorem}

% ---------- 参考文献（BibTeX）----------
% 数据源: refs.bib (76 条, key = klai1..klai76, 顺序与原书 [1]--[76] 一致)
% 样式: unsrt —— 按引用出现顺序编号（即 \cite 首次出现的顺序）
% 说明: 采用 unsrt 以保证编号尽量贴近原文; 若希望按作者字母序可用 plain.
\renewcommand{\bibname}{参考文献}
\bibliographystyle{unsrt}
\nocite{*}
\bibliography{refs}


\end{document}
